\documentclass[11pt]{amsart}
\usepackage{amssymb,amsmath}
\usepackage[margin=1.1in]{geometry}
\usepackage{url}
\sloppy

\newtheorem{theorem}{Theorem}[section]
\newtheorem{lemma}[theorem]{Lemma}
\newtheorem{proposition}[theorem]{Proposition}
\newtheorem{corollary}[theorem]{Corollary}
\theoremstyle{definition}
\newtheorem{definition}[theorem]{Definition}
\theoremstyle{remark}
\newtheorem{remark}[theorem]{Remark}

\newcommand{\cH}{\mathcal H}
\newcommand{\ceil}[1]{\lceil #1\rceil}
\newcommand{\floor}[1]{\lfloor #1\rfloor}

\title[Transversals of $(7,2)$-hypergraphs]{Transversals of $(7,2)$-hypergraphs: a certificate-free bound below $7/8$ and the $3/4$ bound for threshold families}
\author{[Author names to be added]}
\date{Draft, 23 September 2026. Not submitted; not for distribution.}

\begin{document}

\begin{abstract}
Let $f(k,7)$ be the largest transversal number of a $k$-uniform hypergraph in which any at most seven edges can be met by two points. Fon-Der-Flaass, Kostochka and Woodall proved $\ceil{3k/4}\le f(k,7)\le\ceil{7k/8}$. Erd\H{o}s asked whether $f(k,7)=(3/4+o(1))k$; the question is open (Erd\H{o}s Problem \#644). We prove
\[
f(k,7)\le\ceil{173k/200}+10\qquad(k\ge1000),
\]
so $\limsup_k f(k,7)/k\le 0.865<7/8$. The proof is entirely by hand. Its core is a local closing lemma for a \emph{good triple} (three edges with no common point); a maximality argument and two intersection-gap lemmas complete it.
Second, we prove the conjectured bound $\tau\le\frac34k+O(1)$ for a natural class of highly symmetric families. A \emph{threshold family} consists of all $k$-sets that meet some prescribed part $P_i$ of a partition in at least $t_i$ points. With at least three such parts, every threshold family with the $(7,2)$ property satisfies $\tau\le\frac34k+8p$. Here $p$ is the number of parts, and $3/4$ is sharp. The proof is a short convexity argument: feasibility of a single Fano-type construction is convex in the parameters, and it can be checked by hand at the vertices of the parameter domain, all of which are degenerate.
\end{abstract}
\maketitle

\noindent\textit{Status of this draft.} This manuscript was prepared with AI assistance. Several lemmas are re-derivations, in uniform notation and with complete proofs, of lemmas from an unpublished working note (see Remark~\ref{rem:prov}). Nothing here has been refereed externally.

\section{Introduction}

A \emph{transversal} of a family of sets is a set meeting every member; $\tau(\cH)$ denotes the minimum size of a transversal of $\cH$. For integers $p>t\ge1$ and $k\ge2$, say that $\cH$ has \emph{property $(p,t)$} if every subfamily of at most $p$ edges has a transversal of size at most $t$. Erd\H{o}s, Hajnal and Tuza~\cite{EHT} raised the problem of determining the largest value $f(k,p,t)$ of $\tau(\cH)$ over $k$-uniform families $\cH$ with property $(p,t)$ (as reported in~\cite{FKW}). We write $f(k,p)=f(k,p,2)$. (In~\cite{FKW} the property is stated for subfamilies of exactly $p$ edges; for families with at least $p$ edges the two versions coincide, and families with at most six edges have $\tau\le6$, which is irrelevant for the asymptotic questions below.) The same parameter is studied in the notation $h_r(k,\ell)$ by Buci\'c, Kor\'andi and Sudakov~\cite{BKS}, who survey its history.

Erd\H{o}s, Fon-Der-Flaass, Kostochka and Tuza~\cite{EFKT} determined $f(k,p)$ for $3\le p\le6$; in particular $f(k,6)=k$, and for these $p$ the extremal families are the largest complete $k$-uniform hypergraphs with property $(p,2)$ (both statements as reported in~\cite{FKW}). The explicit values $f(k,3)=2k$, $f(k,4)=\floor{3k/2}$ and $f(k,5)=\floor{5k/4}$ are recorded in our sources only at second hand [citation to check]. For $p=7$, Fon-Der-Flaass, Kostochka and Woodall~\cite{FKW} proved
\[
f(k,7)\le\ceil{7k/8}\qquad(k\ge8),
\]
and showed that complete hypergraphs are no longer extremal: for $m\ge10$ (and, by their remark, already for $m\ge4$) the family of all $4m$-subsets of a $(7m+1)$-set that meet a fixed $4m$-set in an odd number of points has property $(7,2)$ and transversal number $3m+1$, whereas every complete $4m$-uniform hypergraph with property $(7,2)$ has transversal number at most $3m$. Complete hypergraphs give $f(k,7)\ge\ceil{3k/4}$ (Proposition~\ref{prop:lower}). Erd\H{o}s asked whether
\[
f(k,7)=\Bigl(\tfrac34+o(1)\Bigr)k ;
\]
the paper~\cite{FKW} describes itself as partially answering a question of~\cite{EFKT}, and the question is listed as Problem~\#644 on the Erd\H{o}s problems website~\cite{EP}, where it is marked open. The best published bounds we know of are therefore $3/4\le\liminf f(k,7)/k\le\limsup f(k,7)/k\le7/8$. We have not been able to consult~\cite{Kos}, which treats property $(p,2)$ for large $p$; the Erd\H{o}s problems page~\cite{EP} lists no improvement for $p=7$ [literature search not exhaustive].

\begin{theorem}\label{thm:main}
For every integer $k\ge1000$,
\[
f(k,7)\le\ceil{173k/200}+10 .
\]
In particular $\limsup_{k\to\infty}f(k,7)/k\le173/200=0.865<7/8$. By Lemma~\ref{lem:pad} the same bound holds for families of rank at most $k$.
\end{theorem}

Our second result concerns families with many symmetries. The extremal example $K^{(k)}_{\lfloor7k/4\rfloor-1}$ is invariant under all permutations of its ground set. Given a partition $V=P_1\sqcup\dots\sqcup P_p$, a set $I$ of \emph{boxes} and integers $1\le t_i\le k$, the \emph{threshold family} $\mathcal T(P;I;t)$ consists of all $k$-sets $E$ with $|E\cap P_i|\ge t_i$ for some $i\in I$. It is invariant under $\prod_i\mathrm{Sym}(P_i)$.

\begin{theorem}[Theorem~\ref{thm:threshold}]\label{thm:intro-threshold}
If a threshold family $\mathcal T$ with at least three boxes, each with $|P_i|\ge t_i+7$, has property $(7,2)$, then $\tau(\mathcal T)\le\frac34k+8p$.
\end{theorem}

The constant $3/4$ cannot be improved, and the three-box hypothesis is exactly what Fano-type constructions need (Remark~\ref{rem:threshold}). The proof in Section~\ref{sec:threshold} is short. Feasibility of one fixed Fano construction is a convex condition on the parameters. The vertices of the relevant parameter polytope turn out to be degenerate configurations, and there the construction can be written down by hand.

\begin{remark}[What is proved by hand]
Every step of the proof of Theorem~\ref{thm:main} is a finite chain of explicit inequalities between linear expressions in the parameters; no computer search, linear-programming certificate or case enumeration by machine is used. The appendix lists independent exact computer checks, which are a sanity supplement only.
\end{remark}

\begin{remark}[Provenance]\label{rem:prov}
The coefficient $173/200$, the constants $h=23/50$ and $\ell=43/200$, and Lemmas~\ref{lem:L18}, \ref{lem:L26}, \ref{lem:L32}, \ref{lem:L31}, \ref{lem:gapext}, \ref{lem:L41} and~\ref{lem:finish} below come from an unpublished, computer-assisted working note (Lemmas 7.18, 7.26, 7.32, 7.31, 7.27, 7.41 and 7.50 there). In that note the bound $\ceil{173k/200}+10$ relied on an exact finite certificate (11940 tetrahedral proof nodes) for the exclusion of pair intersections in $[27k/100,23k/50]$; the note also reports stronger computer-assisted coefficients (down to $6/7$), which have not been independently reviewed and are not used here. The new ingredients of the present paper are Proposition~\ref{prop:local}, which replaces that certificate by an eight-case hand argument, the closed-form budgets of the static templates in Lemmas~\ref{lem:S1} and~\ref{lem:S2} (whose label patterns occur among the machine-enumerated templates of the note), and the organisation of the proof around a maximality threshold at $23k/50$. Lemma~\ref{lem:L18} is a parametrised form of Case~2 of~\cite[Lemma~1]{FKW}.
\end{remark}

\subsection*{Strategy}
Let $r=k$, $\beta=173/200$, and suppose an $r$-uniform family with property $(7,2)$ has $\tau>K:=\ceil{\beta r}+10$. Following~\cite{FKW}, we build bad subfamilies of at most seven edges from a \emph{good triple} (three edges with no common point) by \emph{requests}: since $\tau>K$, every set of at most $K$ points is avoided by some edge. Let $m$ be the largest pair intersection not exceeding $\floor{23r/50}$. A balanced request produces a good triple with intersections $m\ge y\ge z$ and $m+2y\le227r/200$ (the maximality of $m$ gives $y,z\le m$). Proposition~\ref{prop:local} closes every such triple unless $m<43r/200$. Hence there is no pair intersection in $[43r/200,23r/50]$. Lemma~\ref{lem:gapext} extends the gap up to $r/2$, and then the conditional finishing Lemma~\ref{lem:finish} gives $\tau\le\ceil{\beta r}+4$, a contradiction.

Section~\ref{sec:prelim} collects preliminaries, Section~\ref{sec:closing} the closing lemmas for good triples, Section~\ref{sec:local} the local Proposition, Section~\ref{sec:gap} the gap lemmas and Section~\ref{sec:main} the proof of Theorem~\ref{thm:main}. Section~\ref{sec:further}, which is logically independent of the main theorem, records structural lemmas aimed at the $3/4$ question. Appendix~\ref{app:checks} describes computer checks.

\section{Preliminaries}\label{sec:prelim}

A \emph{family} (hypergraph) $\cH$ is a finite collection of finite nonempty sets (edges) on a vertex set $V$; it is \emph{$r$-uniform} if all edges have $r$ points and has \emph{rank at most $k$} if all edges have at most $k$ points. A \emph{bad subfamily} is a set of at most seven edges without a transversal of size at most $2$; thus $\cH$ has property $(7,2)$ if and only if it has no bad subfamily. A \emph{pair intersection} is $|E\cap F|$ for distinct edges $E,F$. For sets $E_1,\dots,E_j$ and a point $v$ let $\sigma(v)=\{i:v\in E_i\}$.

\begin{lemma}[requests]\label{lem:request}
If $\tau(\cH)>T$ and $|D|\le T$, then some edge of $\cH$ is disjoint from $D$.
\end{lemma}
\begin{proof}
Otherwise $D$ is a transversal of size at most $T$.
\end{proof}
We call this \emph{requesting an edge avoiding $D$}; the edge is the \emph{response}. Requests may depend on earlier responses.

\begin{lemma}[closing principle]\label{lem:close}
Let $E_1,\dots,E_j$ be edges with $E_1\cap\dots\cap E_j=\emptyset$, and let $D_1,\dots,D_s$, $j+s\le7$, be sets such that every $2$-set $\{p,q\}$ with $\sigma(p)\cup\sigma(q)=\{1,\dots,j\}$ is contained in some $D_i$. If $H_i$ is an edge disjoint from $D_i$ ($1\le i\le s$), then $\{E_1,\dots,E_j,H_1,\dots,H_s\}$ is a bad subfamily.
\end{lemma}
\begin{proof}
A transversal $\{p,q\}$ of size at most two ($p=q$ allowed) of $E_1,\dots,E_j$ satisfies $\sigma(p)\cup\sigma(q)=\{1,\dots,j\}$, and $p\ne q$ since no point lies in all $E_i$. So $\{p,q\}\subseteq D_i$ for some $i$ and $H_i$ misses both points. The subfamily has at most seven distinct members; coincidences among the edges only make it smaller.
\end{proof}

\begin{lemma}[candidate pairs of a good triple]\label{lem:cand3}
Let $E,F,G$ be edges with $E\cap F\cap G=\emptyset$ (a \emph{good triple}). Put
\[
X=E\cap F,\quad Y=E\cap G,\quad Z=F\cap G,\quad P_E=E\setminus(F\cup G),\quad P_F=F\setminus(E\cup G),\quad P_G=G\setminus(E\cup F).
\]
A $2$-set meets $E$, $F$ and $G$ if and only if it has one point in each of the two sets of one of the six \emph{candidate products}
\[
X\times Y,\quad X\times Z,\quad Y\times Z,\quad X\times P_G,\quad Y\times P_F,\quad Z\times P_E .
\]
If the family is $r$-uniform and $x,y,z$ denote $|X|,|Y|,|Z|$, then $|P_E|=r-x-y$, $|P_F|=r-x-z$, $|P_G|=r-y-z$.
\end{lemma}
\begin{proof}
Points outside $E\cup F\cup G$ have $\sigma=\emptyset$ and would need a partner in all three edges. Every other point has as $\sigma$ a nonempty proper subset of $\{E,F,G\}$; two such sets have union $\{E,F,G\}$ exactly when they are two distinct $2$-sets, or a $2$-set and the complementary singleton. The six cells $X,Y,Z,P_E,P_F,P_G$ are disjoint, which gives the sizes.
\end{proof}

\begin{remark}[relabelling]\label{rem:relabel}
Permuting $E,F,G$ permutes the three cells $X,Y,Z$, and every permutation of the cells arises; each private part stays attached to the cell it is paired with ($X$ with $P_G$, $Y$ with $P_F$, $Z$ with $P_E$). Consequently every lemma of Section~\ref{sec:closing}, stated for cells named $X,Y,Z$ with sizes $x,y,z$, applies to a good triple with the three cells assigned to the roles $X,Y,Z$ in any order. We write, for instance, ``Lemma~\ref{lem:L31} with $(x,y,z):=(z,y,m)$'' for the application in which the cell of size $z$ plays the role $X$, the cell of size $y$ the role $Y$ and the cell of size $m$ the role $Z$.
\end{remark}

\begin{lemma}[four edges]\label{lem:cand4}
Let $E,F,G,H$ be edges with $E\cap F\cap G=\emptyset$, and put $X=E\cap F$, $Y=E\cap G$, $Z=F\cap G$, $A=F\cap H$, $B=G\cap H$, $C=E\cap H$.
\begin{enumerate}
\item If $H$ is disjoint from $X\cup Y\cup Z$, then a $2$-set meets all four edges if and only if it lies in one of $X\times B$, $Y\times A$, $Z\times C$.
\item If $H$ is disjoint from $X\cup Y$, put $Q=Z\cap H$, $A'=A\setminus Q$, $B'=B\setminus Q$. Then a $2$-set meets all four edges if and only if it lies in one of $Q\times E$, $X\times B'$, $Y\times A'$, $(Z\setminus Q)\times C$.
\end{enumerate}
\end{lemma}
\begin{proof}
(1) All triple intersections of the four edges are empty, so every $\sigma$ has size at most two, and two such sets cover the four edges only if they are complementary $2$-sets: $\{E,F\},\{G,H\}$, or $\{E,G\},\{F,H\}$, or $\{F,G\},\{E,H\}$. The corresponding cells are $X$ and $B$, $Y$ and $A$, $Z$ and $C$.
(2) Now the only triple intersection is $Q=F\cap G\cap H$, of type $\{F,G,H\}$; its points pair with exactly the points whose type contains $E$, i.e.\ the points of $E$. Points of type $\{G,H\}$, $\{F,H\}$, $\{F,G\}$, $\{E,H\}$ form $B'$, $A'$, $Z\setminus Q$ and $C$ respectively, and the complementary $2$-type pairs are as in (1). Singleton types need a partner of type of size three, i.e.\ in $Q$, whose partner must contain $E$; this is the product $Q\times E$ again.
\end{proof}

\begin{lemma}[balanced request]\label{lem:balanced}
Let $\cH$ be $r$-uniform, let $E\ne G$ be edges with $|E\cap G|=q$, and let $T_0$ be an integer with $q\le T_0\le r$ and $\tau(\cH)>T_0$, $T_0+q\ge2$. There is an edge $F\notin\{E,G\}$ with $E\cap F\cap G=\emptyset$ and
\[
|E\cap F|,\ |G\cap F|\ \le\ r-\floor{(T_0+q)/2}\ \le\ r-\frac{T_0+q-1}{2}.
\]
\end{lemma}
\begin{proof}
Choose $E'\subseteq E$ and $G'\subseteq G$ containing $E\cap G$ with $|E'|=\floor{(T_0+q)/2}$ and $|G'|=\ceil{(T_0+q)/2}$; this is possible because $q\le T_0$ and $T_0+q\le 2r-1$ (as $q\le r-1$ for distinct edges). Then $E'\cap G'=E\cap G$ and $|E'\cup G'|=T_0$. Request $F$ avoiding $E'\cup G'$ (Lemma~\ref{lem:request}). Then $|E\cap F|\le r-|E'|$, $|G\cap F|\le r-|G'|$, $F$ misses $E\cap G$, and $F\ne E,G$ because $E',G'$ are nonempty.
\end{proof}

\begin{lemma}[padding]\label{lem:pad}
Let $\cH$ have rank at most $k$ and property $(7,2)$. Adding to each edge $E$ a set of $k-|E|$ new vertices private to $E$ gives a $k$-uniform family $\cH'$ with property $(7,2)$ and $\tau(\cH')=\tau(\cH)$.
\end{lemma}
\begin{proof}
Each padded edge contains the original, so transversals of subfamilies of $\cH$ are transversals of the corresponding subfamilies of $\cH'$; this gives property $(7,2)$ and $\tau(\cH')\le\tau(\cH)$. Conversely, in a transversal of $\cH'$ replace each private vertex of an edge $E$ by any point of $E$ (edges are nonempty); the result is a transversal of $\cH$ of no larger size.
\end{proof}

\begin{proposition}[the lower bound]\label{prop:lower}
For every $k\ge2$, the complete $k$-uniform hypergraph on $N=\ceil{7k/4}-1$ points has property $(7,2)$ and transversal number $\ceil{3k/4}$. Hence $f(k,7)\ge\ceil{3k/4}$.
\end{proposition}
\begin{proof}
Seven $k$-sets $A_1,\dots,A_7$ of $[N]$ (repetitions allowed; this also covers subfamilies of fewer than seven edges) have a transversal $\{v,w\}$ ($v=w$ allowed) if and only if no $A_i$ misses both $v$ and $w$. Hence they have no transversal of size at most two if and only if their complements $P_i=[N]\setminus A_i$ cover every point and every pair of $[N]$. Suppose they do. If every point lies in at least three $P_i$ then $\sum|P_i|\ge3N$; otherwise some point $a$ lies in at most two blocks, say only in $P_i$ and $P_j$ ($i=j$ allowed), every other point shares a block with $a$ and hence lies in $P_i\cup P_j$, so $|P_i|+|P_j|\ge N$. Either way some $|P_i|\ge3N/7$. But $|P_i|=N-k<3N/7$ because $N<7k/4$. Finally, a set is a transversal of the complete hypergraph iff its complement has fewer than $k$ points, so $\tau=N-k+1=\ceil{7k/4}-k=\ceil{3k/4}$.
\end{proof}

\section{Closing lemmas for a good triple}\label{sec:closing}

\noindent\textbf{Standing assumptions for this section.} $\cH$ is $r$-uniform, $T$ is an integer, $E,F,G$ is a good triple with the notation of Lemma~\ref{lem:cand3}, $x=|X|$, $y=|Y|$, $z=|Z|$ and $S=x+y+z$. Each lemma concludes that $\cH$ contains a bad subfamily consisting of $E,F,G$ and at most four responses. By Lemma~\ref{lem:close} it suffices, at each final stage, to exhibit requests (sets of at most $T$ points, where $\tau(\cH)>T$) that together contain every $2$-set meeting all edges constructed so far. All subset sizes below are integers whenever the data are integers.

\begin{lemma}[small triple]\label{lem:L18}
Suppose $x,y,z\le m\le r/2$ and
\[
\tau(\cH)>B:=\Bigl\lceil\max\Bigl\{\frac{3r+m}{4},\ \frac{2r+2m}{3}\Bigr\}\Bigr\rceil .
\]
Then $\cH$ contains a bad subfamily.
\end{lemma}
\begin{proof}
The construction is a parametrised form of Case~2 in the proof of~\cite[Lemma~1]{FKW}. Since the hypothesis is symmetric, by Remark~\ref{rem:relabel} we may assume $x\ge y\ge z$. Note $B\le r$ (both terms are at most $r$ as $m\le r/2$) and $B\ge 2m$ (since $(2r+2m)/3\ge2m$). Choose
\begin{itemize}
\item $A_F\subseteq P_F$ with $|A_F|=B-x-y$; this is possible since $0\le B-2m\le B-x-y\le r-x-z=|P_F|$, the last because $B\le r$ and $y\ge z$;
\item $B_0\subseteq P_G$ with $|B_0|=\min(B-x,\ r-y-z)$, and put $C:=G\setminus(Y\cup B_0)=Z\cup(P_G\setminus B_0)$, so $c:=|C|=\max(r-B+x-y,\ z)$;
\item $A_E\subseteq P_E$ with $|A_E|=\min(B-x-c,\ r-x-y)$; here $B-x-c\ge0$ because $B-x-z\ge B-2m\ge0$ and $B-x-(r-B+x-y)\ge 2B-r-2m\ge0$ (as $2B\ge(4r+4m)/3\ge r+2m$).
\end{itemize}
Request edges avoiding
\[
D_1=X\cup Y\cup A_F,\quad D_2=X\cup B_0,\quad D_3=X\cup C\cup A_E,\quad D_4=Y\cup Z\cup(P_E\setminus A_E)\cup(P_F\setminus A_F).
\]
\emph{Sizes.} $|D_1|=B$, $|D_2|\le x+(B-x)=B$, $|D_3|\le x+c+(B-x-c)=B$. For $D_4$,
\[
|D_4|=y+z+\bigl(r-x-y-|A_E|\bigr)+\bigl(r-x-z-(B-x-y)\bigr)=r+2y-B+\max(0,\ r-y-B+c),
\]
and substituting $c$ gives $|D_4|=\max\{r-B+2y,\ 3(r-B)+x,\ 2(r-B)+y+z\}$. Each term is at most $B$: $2B\ge r+2m\ge r+2y$ (as $2B\ge(4r+4m)/3\ge r+2m$), $4B\ge3r+m\ge3r+x$, and $3B\ge2r+2m\ge2r+y+z$. All four requests have at most $B<\tau(\cH)$ points.

\emph{Coverage} (Lemma~\ref{lem:cand3}). $X\times Y\subseteq D_1$. $X\times Z\subseteq D_3$ since $Z\subseteq C$. $X\times P_G$: $P_G=B_0\cup(P_G\setminus B_0)$ with $B_0\subseteq D_2$ and $P_G\setminus B_0\subseteq C\subseteq D_3$. $Y\times Z\subseteq D_4$. $Y\times P_F$: points of $A_F$ are with $Y$ in $D_1$, the rest in $D_4$. $Z\times P_E$: points of $A_E$ are with $Z$ in $D_3$, the rest in $D_4$. Lemma~\ref{lem:close} finishes the proof.
\end{proof}

\begin{lemma}[three response cases]\label{lem:L26}
Suppose
\[
T\ge\max\Bigl\{S,\ r-x+z,\ r-y+z,\ r-x+\frac y2,\ r-y+\frac x2,\ \frac{r+2x+2y+z}{3}\Bigr\}
\]
and $\tau(\cH)>T$. Then $\cH$ contains a bad subfamily.
\end{lemma}
\begin{proof}
Request $H$ avoiding $X\cup Y\cup Z$ ($S\le T$ points). With $A=F\cap H$, $B=G\cap H$, $C=E\cap H$ (sizes $a,b,c$) we have $A\subseteq P_F$, $B\subseteq P_G$, $C\subseteq P_E$, hence
\[
a\le r-x-z,\qquad b\le r-y-z,\qquad c\le r-x-y,\qquad a+b+c\le r,
\]
and by Lemma~\ref{lem:cand4}(1) it remains to cover $X\times B$, $Y\times A$, $Z\times C$ with three requests of at most $T$ points. Note $T-x,T-y\ge0$ since $T\ge S$.

\emph{Case $b\le T-x$.} Use $X\cup B$, $Y\cup Z\cup C\cup A_1$, $Y\cup A_2$, where $A=A_1\sqcup A_2$ with $|A_1|\le T-y-z-c$ and $|A_2|\le T-y$. Such a split exists because $T-y-z-c\ge T-y-z-(r-x-y)=T-(r-x+z)\ge0$ and $(T-y-z-c)+(T-y)\ge a$, i.e.\ $a+z+c\le 2T-2y$, which holds since $a+z+c\le(r-x-z)+z+(r-x-y)=2r-2x-y\le 2T-2y$ by $T\ge r-x+y/2$.

\emph{Case $a\le T-y$.} Symmetric: use $Y\cup A$, $X\cup Z\cup C\cup B_1$, $X\cup B_2$, using $T\ge r-y+z$ and $T\ge r-y+x/2$.

\emph{Case $b>T-x$ and $a>T-y$.} Choose $B_2\subseteq B$ with $|B_2|=T-x$ and $A_3\subseteq A$ with $|A_3|=T-y$, and use
\[
X\cup B_2,\qquad Y\cup A_3,\qquad X\cup Y\cup Z\cup C\cup(A\setminus A_3)\cup(B\setminus B_2).
\]
The first two have exactly $T$ points; the third has $S+c+a-(T-y)+b-(T-x)=a+b+c+2x+2y+z-2T\le r+2x+2y+z-2T\le T$ points. In each case every candidate product lies in one request.
\end{proof}

\begin{lemma}[padding one private part]\label{lem:L32}
Suppose $T\le r$, $\tau(\cH)>T$ and
\[
T\ge\max\Bigl\{S,\ \frac r2+y,\ \frac{r+2x-y+z}{2},\ \frac{r+2x+y+3z}{3}\Bigr\}.
\]
Then $\cH$ contains a bad subfamily.
\end{lemma}
\begin{proof}
Request $H$ avoiding $X\cup Y\cup Z$ together with $T-S$ points of $P_F$; this fits because $0\le T-S\le r-x-z$, the latter being $T\le r+y$. With $A,B,C$ as in Lemma~\ref{lem:L26}, now $a\le r-x-z-(T-S)=r-T+y$, $b\le r-y-z$, $a+b+c\le r$, and it remains to cover $X\times B$, $Y\times A$, $Z\times C$. Note $2(T-x-z)\ge r-y-z\ge b\ge0$ by the third term.

\emph{Case $a\le T-y-z$.} Split $B=B_1\sqcup B_2$ with $|B_i|\le T-x-z$ and start the requests as $X\cup Z\cup B_1$, $X\cup Z\cup B_2$, $Y\cup Z\cup A$, all of size at most $T$. Their total size plus $c$ is at most $2x+y+3z+a+b+c\le r+2x+y+3z\le 3T$, so $C$ (which is disjoint from all three) can be distributed into the residual capacities. Every request contains $Z$, so $Z\times C$ is covered; the first two cover $X\times B$ and the third covers $Y\times A$.

\emph{Case $a>T-y-z$.} Then $b+c\le r-a<r-T+y+z\le 2(T-x-z)$, the last inequality being the fourth term. Split $B\cup C$ into two parts of size at most $T-x-z$ and use $X\cup Z\cup(\text{part}_1)$, $X\cup Z\cup(\text{part}_2)$ and $Y\cup A$; the last has at most $y+r-T+y\le T$ points by the second term. Again every candidate product is covered.
\end{proof}

\begin{lemma}[a surviving triple cell]\label{lem:L31}
Suppose $T\le r$, $\tau(\cH)>T$ and
\[
T\ge\max\Bigl\{x+y,\ \frac r2+x,\ \frac r2+y,\ r+x-y-z,\ r-x+y-z,\ r-\frac S3,\ \frac{3r+S}{5},\ \frac{r+x+y+2z}{3},\ \frac{2r+3z}{4}\Bigr\}.
\]
Then $\cH$ contains a bad subfamily.
\end{lemma}
\begin{proof}
Request $H$ avoiding $X\cup Y\cup Z_0$, where $Z_0\subseteq Z$ has size $\min(z,T-x-y)$. Put $Q=Z\cap H$, $A=(F\cap H)\setminus Q$, $B=(G\cap H)\setminus Q$, $C=E\cap H$, $V=Z\setminus Q$, $D=E\setminus(X\cup Y\cup C)$, with sizes $q,a,b,c,v=z-q,d=r-x-y-c$. Then $A\subseteq P_F$, $B\subseteq P_G$, $C\subseteq P_E$ and
\[
q\le(S-T)_+,\quad a\le r-x-z,\quad b\le r-y-z,\quad c\le r-x-y,\quad q+a+b+c\le r .
\]
By Lemma~\ref{lem:cand4}(2) we must cover $Q\times E$, $X\times B$, $Y\times A$, $V\times C$. We shall use three requests that all contain $Q$ and whose union contains $E$; then $Q\times E$ is covered.

\emph{Two uniform bounds.} $q+x+b\le\max(r+x-y-z,\ r+2x-T)\le T$. Indeed if $S\le T$ then $q=0$ and $b\le r-y-z$; if $S>T$ then $|G\cap(X\cup Y\cup Z_0)|=y+(T-x-y)=T-x$, so $q+b=|G\cap H|\le r-T+x$. Symmetrically $q+y+a\le\max(r-x+y-z,\ r+2y-T)\le T$ (using $|F\cap(X\cup Y\cup Z_0)|=T-y$). Also $z\le T$, since $4T\ge2r+3z\ge5z$.

\emph{Case 0: $c\le T-z$.} Start with $Q\cup X\cup B$, $Q\cup Y\cup A$, $Z\cup C$, each of size at most $T$. Their total plus $d$ equals $r+2q+a+b+z$. Using $a+b\le 2r-x-y-2z$: if $S\le T$ this is at most $3r-S\le3T$ (term $r-S/3$); if $S>T$ it is at most $3r-S+2(S-T)=3r+S-2T\le3T$ (term $(3r+S)/5$). Distribute $D$ (disjoint from all three) into the residual capacities.

In the remaining cases $c>T-z$; put $a_0=r+x+z-2T$ and $b_0=r+y+z-2T$.

\emph{Case 1: $a<a_0$.} Start with $Q\cup X\cup B$, $Y\cup A\cup Z\cup C_2$, $Z\cup C_3$, where $C=C_2\sqcup C_3$, $|C_2|\le T-y-a-z$, $|C_3|\le T-z$. This is possible: $y+a+z< r+x+y+2z-2T\le T$ (term $(r+x+y+2z)/3$), and $a+c< a_0+(r-x-y)=2r+z-y-2T\le 2T-2z-y$ (term $(2r+3z)/4$), i.e.\ the two capacities sum to more than $c$. The total plus $d$ is $r+q+a+b+2z\le 2r+2z-c<2r+3z-T\le3T$. Distribute $D$.

\emph{Case 2: $b<b_0$.} Symmetric to Case 1, exchanging $(X,B,x,b)$ with $(Y,A,y,a)$.

\emph{Case 3: $a\ge a_0$ and $b\ge b_0$.} Put $u=c+z-T$; then $0<u\le c$ since $z\le T$. Choose $C_{12}\subseteq C$ with $|C_{12}|=u$, put $C_3=C\setminus C_{12}$ (so $|C_3|=T-z$), and choose an integer $t$ with
\[
\max(0,\ y+a+z+u-T)\ \le\ t\ \le\ \min(v,\ T-q-x-b-u).
\]
The interval is nonempty: $T-q-x-b-u=2T-x-z-(q+b+c)\ge0$ because $q+b+c\le r-a\le r-a_0$; $y+a+z+u-T\le v$ is $q+a+c\le 2T-y-z$, true because $q+a+c\le r-b\le r-b_0$; and $y+a+z+u-T\le T-q-x-b-u$ is $q+a+b+2c+x+y+3z\le4T$, which follows from $q+a+b+c\le r$, $c\le r-x-y$ and $2r+3z\le4T$. Split $V=V_{13}\sqcup V_{23}$ with $|V_{13}|=t$ and start with
\[
Q\cup X\cup B\cup C_{12}\cup V_{13},\qquad Q\cup Y\cup A\cup C_{12}\cup V_{23},\qquad Z\cup C_3 .
\]
The choice of $t$ makes the first two of size at most $T$ (using $q+v=z$), the third has exactly $T$ points, and the total plus $d$ is $r+q+a+b+c+3z-T\le 2r+3z-T\le 3T$. Distribute $D$.

In every case $X\times B$ and $Y\times A$ are covered by the first two requests, and $V\times C$ is covered since each part of $C$ lies in a request together with all of $V$ (in Case 3, $C_{12}$ lies in both of the first two requests, which together contain $V$, and $C_3$ lies with $Z\supseteq V$ in the third). All three requests contain $Q$ and their union contains $X\cup Y\cup C\cup D=E$.
\end{proof}

\begin{lemma}[static template S1]\label{lem:S1}
Suppose $\tau(\cH)>T$ and there are integers $x_1,y_1,z_1$ with $0\le x_1\le x$, $0\le y_1\le y$, $0\le z_1\le z$ and
\[
x_1+y_1+z_1\le T,\qquad y_1+z_1\ge r+x-T,\qquad x_1+z_1\ge r+y-T,\qquad x_1+y_1\ge r+z-T.
\]
Then $\cH$ contains a bad subfamily.
\end{lemma}
\begin{proof}
Split $X=X_1\sqcup X_2$, $Y=Y_1\sqcup Y_2$, $Z=Z_1\sqcup Z_2$ with $|X_1|=x_1$, $|Y_1|=y_1$, $|Z_1|=z_1$, and request (all four at once)
\[
R_0=X_1\cup Y_1\cup Z_1,\quad R_X=X\cup P_G\cup Y_2\cup Z_2,\quad R_Y=Y\cup P_F\cup X_2\cup Z_2,\quad R_Z=Z\cup P_E\cup X_2\cup Y_2 .
\]
Their sizes are $x_1+y_1+z_1$, $r+x-y_1-z_1$, $r+y-x_1-z_1$, $r+z-x_1-y_1$, all at most $T$. Coverage: $X\times P_G\subseteq R_X$, $Y\times P_F\subseteq R_Y$, $Z\times P_E\subseteq R_Z$. For $(p,q)\in X\times Y$: if $q\in Y_2$ then $\{p,q\}\subseteq R_X$; if $p\in X_2$ then $\{p,q\}\subseteq R_Y$; otherwise $(p,q)\in X_1\times Y_1\subseteq R_0$. The products $X\times Z$ (via $R_X$, $R_Z$, $R_0$) and $Y\times Z$ (via $R_Y$, $R_Z$, $R_0$) are handled identically.
\end{proof}

\begin{lemma}[static template S2]\label{lem:S2}
Suppose $\tau(\cH)>T$, and put $P=\max(0,r+y-x-z-T)$ and $Q=\max(0,r+x-y-z-T)$. If $x\le T$, $y\le T$ and
\[
\text{(i) } T\ge r-x+z+P+Q,\qquad \text{(ii) } T\ge y+z+Q,\qquad \text{(iii) } 2T\ge r+y+2z+P+2Q,
\]
then $\cH$ contains a bad subfamily.
\end{lemma}
\begin{proof}
Since $y\le T$ we have $P\le r-x-z=|P_F|$, and since $x\le T$ we have $Q\le r-y-z=|P_G|$. Split $P_F=V_1\sqcup V_2$ with $|V_1|=P$ and $P_G=W_0\sqcup W_{13}$ with $|W_{13}|=Q$. Put $U_1=T-r+x-z-P-Q$ and $L=x+y+z+Q-T$. The integer interval $[\max(0,L),\min(x,U_1)]$ is nonempty: $0\le U_1$ is (i), $L\le x$ is (ii), $L\le U_1$ is (iii), and $0\le x$. Choose $x_{01}$ in it and split $X=X_{01}\sqcup X_{03}$ with $|X_{01}|=x_{01}$. Request
\[
R_0=X\cup W_0,\quad R_1=X_{01}\cup Y\cup Z\cup P_E\cup V_1\cup W_{13},\quad R_2=Y\cup V_2,\quad R_3=X_{03}\cup Y\cup Z\cup W_{13}.
\]
Sizes: $|R_0|=x+r-y-z-Q\le T$ by the definition of $Q$; $|R_2|=y+r-x-z-P\le T$ by the definition of $P$; $|R_1|=x_{01}+r-x+z+P+Q\le T$ as $x_{01}\le U_1$; $|R_3|=x-x_{01}+y+z+Q\le T$ as $x_{01}\ge L$. Coverage: $X\times P_G$: $X\times W_0\subseteq R_0$, $X_{01}\times W_{13}\subseteq R_1$, $X_{03}\times W_{13}\subseteq R_3$. $Y\times P_F$: $Y\times V_1\subseteq R_1$, $Y\times V_2\subseteq R_2$. $Z\times P_E\subseteq R_1$. $X\times Y$ and $X\times Z$: $X_{01}$ with $Y\cup Z$ in $R_1$, $X_{03}$ with $Y\cup Z$ in $R_3$. $Y\times Z\subseteq R_1$.
\end{proof}

\begin{lemma}[scaling and monotonicity]\label{lem:scaling}
Each hypothesis of Lemmas~\ref{lem:L26}, \ref{lem:L32}, \ref{lem:L31} and~\ref{lem:S2}, other than $T\le r$, has the form $g(r,x,y,z,T)\ge0$ with $g$ positively homogeneous of degree one and nondecreasing in $T$. Consequently, if these hypotheses hold for the real numbers $(1,x/r,y/r,z/r,\beta)$, then they hold for $(r,x,y,z,T)$ for every real $T\ge\beta r$.
\end{lemma}
\begin{proof}
For Lemmas~\ref{lem:L26}, \ref{lem:L32}, \ref{lem:L31} each hypothesis is $T-\ell(r,x,y,z)\ge0$ or $cT-\ell\ge0$ with $c>0$ and $\ell$ linear. For Lemma~\ref{lem:S2}, $P$ and $Q$ are of the form $\max(0,\ell-T)$, which is positively homogeneous and nonincreasing in $T$; each of $T-x$, $T-y$, $T-(r-x+z)-P-Q$, $T-y-z-Q$, $2T-(r+y+2z)-P-2Q$ is then homogeneous and nondecreasing in $T$. Homogeneity transfers the inequalities from $(1,x/r,y/r,z/r,\beta)$ to $(r,x,y,z,\beta r)$, and monotonicity to every $T\ge\beta r$.
\end{proof}

\section{The local closing proposition}\label{sec:local}

Throughout this section and the next, $\beta=173/200$ and $e:=1-\beta=27/200$.

\begin{proposition}[local closing lemma]\label{prop:local}
Let $r$ and $T$ be integers with $\beta r+3\le T\le r$, and let $\cH$ be an $r$-uniform family with $\tau(\cH)>T$. Let $E,F,G$ be a good triple whose three pair intersections have sizes $m\ge y\ge z$, where
\[
m\le\frac{23r}{50}\qquad\text{and}\qquad m+2y\le\frac{227r}{200}=(2-\beta)r .
\]
Then $\cH$ contains a bad subfamily. No hypothesis on other pair intersections of $\cH$ is needed.
\end{proposition}

\begin{proof}
Write $\bar m=m/r$, $\bar y=y/r$, $\bar z=z/r$. The case distinctions below refer to these normalised values; to lighten notation we drop the bars, so in each case $m,y,z$ are real numbers with
\[
0\le z\le y\le m\le\tfrac{92}{200},\qquad m+2y\le\tfrac{227}{200}. \tag{R}
\]
In cases (b) and (c) we verify the hypotheses of the lemma used at $r=1$, $T=\beta$; by Lemma~\ref{lem:scaling} they then hold for the actual integers $r$ and $T\ge\beta r$ (and $T\le r$ is given). Case (a) and the rounding in case (c4) are treated directly at the integer level.

\medskip\noindent\textbf{(a) $m\le119/400$.} Here we return to the unnormalised size $m\le\tfrac{119}{400}r$. Apply Lemma~\ref{lem:L18} with this $m$ (all three intersections are at most $m\le r/2$). Its budget satisfies $(2r+2m)/3\le(2+2\cdot\tfrac{119}{400})r/3=\beta r$ and $(3r+m)/4\le(3+\tfrac{119}{400})r/4=\tfrac{1319}{1600}r<\beta r$, so $B\le\ceil{\beta r}\le T<\tau(\cH)$. (This case uses its own budget $B$, not $T$.)

\medskip\noindent\textbf{(b) $m>119/400$ and $y\le73/200$.}

\emph{(b2) $y-z>m-e$.} Lemma~\ref{lem:L32} with $(x,y,z):=(m,y,z)$. Its conditions are: $S\le\beta$: from $z<y-m+e$ we get $S<2y+e\le\tfrac{146+27}{200}=\beta$. $\tfrac12+y\le\beta$: $y\le\tfrac{73}{200}$. $\tfrac{1+2m-y+z}{2}\le\beta$, i.e.\ $2m-y+z\le\tfrac{146}{200}$: $2m-y+z<m+e\le\tfrac{92+27}{200}$. $\tfrac{1+2m+y+3z}{3}\le\beta$, i.e.\ $2m+y+3z\le\tfrac{319}{200}$: $2m+y+3z<4y-m+3e\le3y+\tfrac{81}{200}\le\tfrac{219+81}{200}$.

\emph{(b3) $y-z\le m-e$ and $S<81/200$.} Lemma~\ref{lem:L32} with $(x,y,z):=(z,y,m)$. Conditions: $S<\tfrac{81}{200}\le\beta$; $\tfrac12+y\le\beta$; $2z-y+m\le z+m\le S<\tfrac{146}{200}$; $2z+y+3m\le3S<\tfrac{243}{200}\le\tfrac{319}{200}$.

\emph{(b1) $y-z\le m-e$ and $S\ge81/200$.} Lemma~\ref{lem:L31} with $(x,y,z):=(z,y,m)$, so its nine conditions read (with $S=m+y+z$):
\begin{enumerate}
\item $z+y\le\beta$: $z+y\le2y\le\tfrac{146}{200}$.
\item $\tfrac12+z\le\beta$ and (3) $\tfrac12+y\le\beta$: from $z\le y\le\tfrac{73}{200}$.
\setcounter{enumi}{3}
\item $1+z-y-m\le\beta$, i.e.\ $y+m-z\ge e$: $y+m-z\ge m>\tfrac{119}{400}>e$.
\item $1-z+y-m\le\beta$, i.e.\ $y-z\le m-e$: the case hypothesis.
\item $1-S/3\le\beta$, i.e.\ $S\ge3e=\tfrac{81}{200}$: the case hypothesis.
\item $\tfrac{3+S}{5}\le\beta$, i.e.\ $S\le5\beta-3=\tfrac{265}{200}$: $S\le m+2y\le\tfrac{227}{200}$.
\item $\tfrac{1+z+y+2m}{3}\le\beta$, i.e.\ $z+y+2m\le\tfrac{319}{200}$: $z+y\le2y\le\tfrac{227}{200}-m$, so $z+y+2m\le\tfrac{227}{200}+m\le\tfrac{227+92}{200}$. (This is where $m\le23/50$ is used; it is tight at $m=23/50$, $z=y$, $m+2y=227/200$.)
\item $\tfrac{2+3m}{4}\le\beta$, i.e.\ $m\le\tfrac{292}{600}$: $m\le\tfrac{92}{200}$.
\end{enumerate}
Cases (b1), (b2), (b3) exhaust (b).

\medskip\noindent\textbf{(c) $y>73/200$} (then $m\ge y>119/400$). Put $u=m+y$ and $d=m-y\ge0$. From (R) and $m\ge y>\tfrac{73}{200}$:
\begin{equation*}\tag{C}
\begin{gathered}
m\le\tfrac{227}{200}-2y<\tfrac{81}{200},\qquad \tfrac{146}{200}<u<\tfrac{154}{200},\qquad d<\tfrac{8}{200},\\
2m-y\le\tfrac{454}{200}-5y<\tfrac{89}{200},\qquad 2m+3y\le\tfrac{454}{200}-y<\tfrac{381}{200},\qquad 2m+y\le\tfrac{454}{200}-3y<\tfrac{235}{200}.
\end{gathered}
\end{equation*}
(For the last three use $m\le\tfrac{227}{200}-2y$ and $y>\tfrac{73}{200}$.)

\emph{(c1) $z\le\tfrac{319}{200}-2u$.} Lemma~\ref{lem:L26} with $(x,y,z):=(m,y,z)$. Conditions: $\tfrac{1+2u+z}{3}\le\beta$ is the case hypothesis. $S=u+z\le\tfrac{319}{200}-u<\tfrac{173}{200}$. $1-y+z\le\beta$, i.e.\ $z\le y-e$: from $z\le\tfrac{319}{200}-2m-2y$ it suffices that $2m+3y\ge\tfrac{346}{200}$, and indeed $2m+3y>5\cdot\tfrac{73}{200}$; then also $1-m+z\le1-y+z\le\beta$. $1-m+y/2\le\beta$, i.e.\ $m-y/2\ge e$: $m-y/2\ge y/2>\tfrac{36.5}{200}$. $1-y+m/2\le\beta$, i.e.\ $y-m/2\ge e$: $y-m/2>\tfrac{73-40.5}{200}=\tfrac{32.5}{200}$.

\emph{(c2) $z>\tfrac{319}{200}-2u$ and $z<e-d$.} Lemma~\ref{lem:S2} with $(x,y,z):=(y,m,z)$ (the $y$-cell plays the role $X$, the $m$-cell the role $Y$). At $r=1$, $T=\beta$: $P=\max(0,e+d-z)=e+d-z>0$ (since $z<e-d\le e+d$) and $Q=\max(0,e-d-z)=e-d-z>0$. The conditions $X=y\le\beta$ and $Y=m\le\beta$ hold by (C). Then
\begin{itemize}
\item (i) $\beta\ge1-y+z+P+Q=1-y-z+2e$, i.e.\ $z\ge\tfrac{81}{200}-y$;
\item (ii) $\beta\ge m+z+Q=m-d+e=y+e$, i.e.\ $y\le\tfrac{146}{200}$, which holds by (C);
\item (iii) $2\beta\ge1+m+2z+P+2Q=1+m-d-z+3e=1+y-z+3e$, i.e.\ $z\ge y-\tfrac{65}{200}$.
\end{itemize}
Both lower bounds follow from $z>\tfrac{319}{200}-2u$: indeed $\tfrac{319}{200}-2u\ge\tfrac{81}{200}-y$ is $2m+y\le\tfrac{238}{200}$, and $\tfrac{319}{200}-2u\ge y-\tfrac{65}{200}$ is $2m+3y\le\tfrac{384}{200}$; both hold by (C).

\emph{(c3) $z>\tfrac{319}{200}-2u$, $z\ge e-d$ and $z\le\tfrac12(\tfrac{146}{200}-y)$.} Lemma~\ref{lem:S2} with $(x,y,z):=(m,y,z)$. Now $P=\max(0,e-d-z)=0$ and $Q=\max(0,e+d-z)$; $m,y\le\beta$ by (C). The conditions become (i) $1-m+z+Q\le\beta$, (ii) $y+z+Q\le\beta$, (iii) $1+y+2z+2Q\le2\beta$.
\begin{itemize}
\item If $Q=0$ ($z\ge e+d$): (i) is $z\le m-e$, which follows from $z\le\tfrac{73}{200}-\tfrac y2$ and $2m+y\ge\tfrac{219}{200}>1$ (so $\tfrac{73}{200}-\tfrac y2\le m-e$); (ii) $y+z\le\tfrac{73}{200}+\tfrac y2<\tfrac{73+40.5}{200}<\beta$; (iii) is $y+2z\le\tfrac{146}{200}$, the case hypothesis.
\item If $Q>0$ ($z<e+d$): (i) becomes $1+e-y\le\beta$, i.e.\ $y\ge2e=\tfrac{54}{200}$; (ii) becomes $m+e\le\beta$, i.e.\ $m\le\tfrac{146}{200}$; (iii) becomes $1+y+2d+2e\le2\beta$, i.e.\ $2m-y\le\tfrac{92}{200}$, true by (C).
\end{itemize}

\emph{(c4) $z>\tfrac12(\tfrac{146}{200}-y)$} (and $z>\tfrac{319}{200}-2u$, $z\ge e-d$). First,
\[
z\ \ge\ \max\Bigl(e+d,\ u-\tfrac{119}{200}\Bigr). \tag{$*$}
\]
Indeed $\tfrac12(\tfrac{146}{200}-y)\ge e+d$ is equivalent to $2m-y\le\tfrac{92}{200}$, and $\tfrac12(\tfrac{146}{200}-y)\ge u-\tfrac{119}{200}$ is equivalent to $2m+3y\le\tfrac{384}{200}$; both hold by (C). We apply Lemma~\ref{lem:S1} with $(x,y,z):=(m,y,z)$ and splits $(m_1,y_1,z_1)$ chosen as follows (normalised; the requirements at $r=1$, $T=\beta$ are $0\le m_1\le m$, $0\le y_1\le y$, $0\le z_1\le z$, $m_1+y_1+z_1\le\beta$, $y_1+z_1\ge e+m$, $m_1+z_1\ge e+y$, $m_1+y_1\ge e+z$).
\begin{itemize}
\item \emph{Split A, if $3z\ge e+u$:} $m_1=\tfrac12(e+y+z-m)$, $y_1=\tfrac12(e+m+z-y)$, $z_1=\tfrac12(e+m+y-z)$. The three pair sums equal $e+m$, $e+y$, $e+z$ exactly; the total is $\tfrac12(3e+S)\le\beta$ because $S\le m+2y\le\tfrac{227}{200}\le2\beta-3e=\tfrac{265}{200}$. Bounds: $m_1\ge0$ as $z\ge e+d\ge d-e$; $m_1\le m$ iff $z\le3m-y-e$, true as $z\le y$ and $2y+e\le3m$ (since $3m-2y\ge y>e$); $y_1\ge0$ clearly; $y_1\le y$ iff $z\le3y-m-e$, true as $z\le y$ and $m+e\le2y$ (since $d<\tfrac{8}{200}<y-e$); $z_1\ge0$ as $z\le y\le m$; $z_1\le z$ iff $3z\ge e+u$.
\item \emph{Split B, if $3z<e+u$:} $z_1=z$, $y_1=e+m-z$, $m_1=e+y-z$. Pair sums: $y_1+z_1=e+m$, $m_1+z_1=e+y$, $m_1+y_1=2e+u-2z\ge e+z$ iff $3z\le e+u$. Total $2e+u-z\le\beta$ iff $z\ge u-\tfrac{119}{200}$, which is $(*)$. Bounds: $y_1\ge0$ as $z\le m$; $y_1\le y$ iff $z\ge e+d$, which is $(*)$; $m_1\ge0$ as $z\le y$; $m_1\le m$ iff $z\ge e-d$.
\end{itemize}
\emph{Integer rounding.} At the integer level take $\hat m_1=\ceil{m_1r}$, $\hat y_1=\ceil{y_1r}$, $\hat z_1=\ceil{z_1r}$, where $m_1,y_1,z_1$ are computed from the normalised cell sizes. Since the cell sizes $m,y,z$ (unnormalised) are integers and $m_1r\le m$ etc., the rounded splits still satisfy $0\le\hat m_1\le m$ etc.; the three lower bounds on pair sums, e.g.\ $\hat y_1+\hat z_1\ge (e+\bar m)r=r+m-\beta r\ge r+m-T$, survive rounding up; and $\hat m_1+\hat y_1+\hat z_1<(m_1+y_1+z_1)r+3\le\beta r+3\le T$. Hence Lemma~\ref{lem:S1} applies with the integer budget $T$. This is the only place where the additive $3$ in $T\ge\beta r+3$ is used.

\medskip\noindent\emph{Exhaustiveness.} The cases are defined by an if-then-else chain: (a); otherwise (b) if $y\le\tfrac{73}{200}$, split into (b2), then (b3), then (b1); otherwise (c), split into (c1), then (c2), then (c3), then (c4). Each later case is the complement of the earlier ones within (c), which is how the inequalities $z>\tfrac{319}{200}-2u$ and $z\ge e-d$ were used in (c2)--(c4). In every case a lemma of Section~\ref{sec:closing} produces a bad subfamily.
\end{proof}

\begin{remark}
The proof never uses property $(7,2)$ or any information about pair intersections other than those of $E,F,G$; it only uses $\tau(\cH)>T$. The condition $m+2y\le(2-\beta)r$ is what a balanced request of size at least $\beta r+1$ guarantees (Lemma~\ref{lem:balanced}), while the bounds $y,z\le m$ on the other two intersections are supplied by the maximality of $m$ in Section~\ref{sec:main}. Neither intersectingness nor any bound on $\nu(\cH)$ is assumed.
\end{remark}

\section{Two lemmas on intersection gaps}\label{sec:gap}

\begin{lemma}[extending a gap to $r/2$]\label{lem:gapext}
Let $h=23/50$, $\ell=43/200$, $r\ge1000$, and let $\cH$ be an $r$-uniform family with $\tau(\cH)>\ceil{\beta r}+10$ and no pair intersection in $[\ell r,hr]$. Then $\cH$ contains a bad subfamily or has no pair intersection in $[hr,r/2]$.
\end{lemma}
\begin{proof}
Put $T=\ceil{\beta r}+4\le r$ and $h_0=\floor{hr}$. Suppose $E,F$ are edges with $X=E\cap F$ and $x=|X|\in[hr,r/2]$. Lemma~\ref{lem:balanced} with $T_0=\ceil{\beta r}$ gives an edge $G\ne E,F$ with $E\cap F\cap G=\emptyset$ and
\[
|E\cap G|,|F\cap G|\le r-\frac{\beta r+x-1}{2}\le\Bigl(1-\frac{\beta+h}{2}\Bigr)r+\frac12=\frac{67.5}{200}r+\frac12<hr .
\]
By the gap, $Y=E\cap G$ and $Z=F\cap G$ have sizes $y,z<\ell r$. Put $S=x+y+z$. Recall $|P_E|=r-x-y$, $|P_F|=r-x-z$, $|P_G|=r-y-z$.

\emph{Case $S\le T$.} Let $p=(r-x-y-h_0)_+$, $t=(r-x-z-h_0)_+$, choose $P_1\subseteq P_E$, $P_2\subseteq P_F$ of sizes $p,t$, and let $D=X\cup Y\cup Z\cup P_1\cup P_2\cup W$, where $W\subseteq P_G$ is chosen so that $|D|=T$. Here $C:=|X\cup Y\cup Z\cup P_1\cup P_2|=x+\max(y,r-x-h_0)+\max(z,r-x-h_0)$, and the four possible values are $S\le T$, $r-h_0+y$ and $r-h_0+z$ (both $<(1-h+\ell)r+1=\tfrac{151}{200}r+1<T$), and $2r-x-2h_0\le(2-3h)r+2=0.62r+2<T$. So $C\le T$, and $|W|=T-C\le T-S\le r-y-z$ (as $T\le r+x$). Request $H$ avoiding $D$. Then $E\cap H\subseteq P_E\setminus P_1$ has at most $\min(r-x-y,h_0)\le hr$ points, and $H\ne E$ (as $H$ misses $X\ne\emptyset$); by the gap $|E\cap H|<\ell r$, and likewise $|F\cap H|<\ell r$. Further $b:=|G\cap H|\le r-y-z-|W|=r-T+x+p+t$. All triple intersections of $E,F,G,H$ are empty, so by Lemma~\ref{lem:cand4}(1) it suffices to cover $X\times B$, $Y\times(F\cap H)$, $Z\times(E\cap H)$, where $B=G\cap H$. Split $B=B_1\sqcup B_2$ with $|B_i|\le\ceil{b/2}$ and request $X\cup B_1$, $X\cup B_2$, $Y\cup Z\cup(E\cap H)\cup(F\cap H)$. The last has fewer than $4\ell r=0.86r<T$ points. For the first two, $x+\ceil{b/2}\le\tfrac12(r+3x+p+t-T+1)$, and expanding the positive parts, $r+3x+p+t$ is at most the maximum of
\[
r+3x,\quad 2r+2x-y-h_0,\quad 2r+2x-z-h_0,\quad 3r+x-y-z-2h_0,
\]
which for $x\le r/2$, $y,z\ge0$, $h_0>hr-1$ is at most $\max(2.5r,\ 2.54r+1,\ 2.58r+2)=2.58r+2$. Since $3T\ge3\beta r+12=2.595r+12$, we get $x+\ceil{b/2}\le T$. By Lemma~\ref{lem:close}, $E,F,G,H$ and the three responses form a bad subfamily.

\emph{Case $S>T$.} Since $x+y<(\tfrac12+\ell)r=0.715r<T$, request $H$ avoiding $X\cup Y\cup Z_0$ with $Z_0\subseteq Z$, $|Z_0|=T-x-y\le z$. With $Q=Z\cap H$, $A=(F\cap H)\setminus Q$, $B=(G\cap H)\setminus Q$, $C=E\cap H$ (sizes $q,a,b,c$) we have $q\le S-T$, $b\le r-y-z$, and
\[
c\le r-x-y<r-T+z<(e+\ell)r-4<hr,\qquad q+a=|F\cap H|\le r-(T-y)<(e+\ell)r-4<hr ,
\]
using $x+y>T-z$, $|F\cap(X\cup Y\cup Z_0)|=T-y$, $y,z<\ell r$ and $e+\ell=0.35<h$. As $H\ne E,F$ (it misses $X$), the gap gives $c<\ell r$ and $q+a<\ell r$. By Lemma~\ref{lem:cand4}(2) we must cover $Q\times E$, $X\times B$, $Y\times A$ and $(Z\setminus Q)\times C$. Start with
\[
R_1=X\cup Q\cup B_1,\qquad R_2=X\cup Q\cup B_2,\qquad R_3=Y\cup Z\cup A\cup C,
\]
where $B=B_1\sqcup B_2$, $|B_i|\le\ceil{b/2}$. Using $q\le x+y+z-T$, $x\le r/2$ and $y+z<2\ell r$,
\[
|R_i|\le x+q+\frac{b+1}{2}\le 2x+\frac{y+z}{2}+\frac r2-T+\frac12<\Bigl(\frac32+\ell-\beta\Bigr)r-\frac72=0.85r-3.5<T\quad(i=1,2),
\]
and $|R_3|<4\ell r<T$. The set $D:=E\setminus(X\cup Y\cup C)$ is disjoint from $R_1,R_2,R_3$, and the total of the three sizes plus $|D|=r-x-y-c$ is
\[
r+x+z+2q+a+b\le 2r+x-y+q+(q+a)<2r+2x+z-T+\ell r\le(3+2\ell)r-T=3.43r-T\le3T,
\]
because $4T\ge3.46r+16$. Distribute $D$ into the residual capacities. Now every request contains $Q$ (as $Q\subseteq Z$) and their union contains $X\cup Y\cup C\cup D=E$, so $Q\times E$ is covered; $R_1,R_2$ cover $X\times B$ and $R_3$ covers $Y\times A$ and $(Z\setminus Q)\times C$. Lemma~\ref{lem:close} gives a bad subfamily of seven edges.
\end{proof}

\begin{lemma}[two large pair cells]\label{lem:L41}
Let $\cH$ be $r$-uniform and suppose every pair intersection is at most $m$ or greater than $r/2$, where $m\le r/4$. Let $E,F,G,H$ be edges with all four triple intersections empty, $X=E\cap F$, $B=G\cap H$ with $x=|X|>r/2$, $b=|B|>r/2$, and suppose the other four pair cells $Y=E\cap G$, $Z=F\cap G$, $A=F\cap H$, $C=E\cap H$ each have at most $m$ points. If $T\le r$ is an integer with $\tau(\cH)>T$ and
\[
T\ge\ceil{r/2}+2m,\qquad T\ge x+m,\qquad 3T\ge2x+b+\ceil{r/2}+2m,
\]
then $\cH$ contains a bad subfamily.
\end{lemma}
\begin{proof}
Put $U=Y\cup Z\cup A\cup C$ (a disjoint union), $s=|Y|+|Z|\le2m$, $t=|A|+|C|\le2m$, and
\[
p=\ceil{r/2}-\min(s,t),\qquad q=T-\ceil{r/2}-\max(s,t).
\]
Then $p\ge0$ (as $2m\le r/2$), $q\ge0$ (first condition), $p\le\ceil{r/2}\le b$, and $q\le T-\ceil{r/2}\le\floor{r/2}<x$. Choose $B_0\subseteq B$, $X_0\subseteq X$ of sizes $p,q$ and request $I$ avoiding $U\cup B_0\cup X_0$, a set of exactly $T$ points. Since $G\setminus(Y\cup Z\cup B_0)$ has $r-s-p=\floor{r/2}+\min(s,t)-s\le\floor{r/2}$ points, $|G\cap I|\le r/2$; likewise $|H\cap I|\le r-t-p\le\floor{r/2}$. Hence $I\ne G,H$ and, by the gap, $|G\cap I|,|H\cap I|\le m$. In particular $d:=|B\cap I|\le m$, and $x_I:=|X\cap I|\le x-q$.

By Lemma~\ref{lem:cand4}(1) the $2$-sets meeting $E,F,G,H$ lie in $X\times B$, $Y\times A$ or $Z\times C$; the last two products lie in $U$ and are missed by $I$. Choose $B_1\subseteq B$ with $B\cap I\subseteq B_1$ and $|B_1|=\min(b,T-x)$ (possible since $T-x\ge m\ge d$), and request edges avoiding
\[
X\cup B_1\qquad\text{and}\qquad (X\cap I)\cup(B\setminus B_1).
\]
The first has at most $T$ points. The second has $x_I+\max(0,b-T+x)$ points; if $b\le T-x$ this is $x_I\le x\le T$, and otherwise it is at most $2x+b-q-T\le 2x+b-T+\ceil{r/2}+2m-T\le T$ by the third condition. Now let $(v,w)\in X\times B$. If $w\in B_1$, the first response misses both points. Otherwise $w\notin I$ (as $B\cap I\subseteq B_1$), so if $\{v,w\}$ meets $I$ then $v\in X\cap I$ and $\{v,w\}$ lies in the second request. So the seven edges $E,F,G,H,I$ and the two responses form a bad subfamily.
\end{proof}

\begin{lemma}[conditional finishing lemma]\label{lem:finish}
Let $5/6\le\beta'<1$, let $r$ be an integer with $T:=\ceil{\beta' r}+4\le r$, and let $\cH$ be an $r$-uniform family with property $(7,2)$ in which every pair intersection is at most $\tfrac{3\beta'-2}{2}r$ or greater than $r/2$. Then $\tau(\cH)\le T$.
\end{lemma}
\begin{proof}
Suppose $\tau(\cH)>T$. Requesting an edge avoiding $T$ points of some edge gives a pair intersection at most $r-T<r/2$, so the largest pair intersection $m$ not exceeding $r/2$ exists, and $m\le\tfrac{3\beta'-2}{2}r\le r/2$. Take $E,G$ with $|E\cap G|=m$ and apply Lemma~\ref{lem:balanced} with $T_0=T$: we obtain $F$ with $E\cap F\cap G=\emptyset$ and $|E\cap F|,|G\cap F|\le r-\floor{(T+m)/2}$.

\emph{If both are at most $r/2$}, they are at most $m$ by maximality, and Lemma~\ref{lem:L18} applies to $E,F,G$: its budget is at most $\ceil{\beta' r}\le T$, because $\tfrac{2r+2m}{3}\le\beta' r$ is equivalent to $m\le\tfrac{3\beta'-2}{2}r$, and $\tfrac{3r+m}{4}\le\tfrac{4+3\beta'}{8}r\le\beta' r$ as $\beta'\ge4/5$. This contradicts property $(7,2)$.

\emph{Otherwise} one of them exceeds $r/2$; exchanging the names $E,G$ we may assume $X:=E\cap F$ has $x>r/2$, and then $Z:=F\cap G$ has $z\le r-x<r/2$, so $z\le m$; put $Y:=E\cap G$ ($|Y|=m$). We have $x\le r-\floor{(T+m)/2}\le r-\tfrac{T+m}{2}+\tfrac12$. Moreover $m\le r-T$: otherwise $T+m\ge r+1$, $\floor{(T+m)/2}\ge\ceil{r/2}$, and $x\le\floor{r/2}$, a contradiction. Hence $m\le r/6-4<r/4$ and
\[
x+2m\le r-\tfrac T2+\tfrac{3m}{2}+\tfrac12\le\tfrac{5r-4T+1}{2}<T .
\]
Request $H$ avoiding $X\cup Y\cup Z$ together with $T-x-m-z$ points of $P_G$; this number is positive (as $x+m+z\le x+2m<T$) and at most $r-m-z=|P_G|$ (as $T\le r+x$). In particular $H\ne G$, and $H\ne E,F$ as $H$ misses $X$. Then $E\cap H\subseteq P_E$ and $F\cap H\subseteq P_F$ have fewer than $r/2$ points and hence at most $m$ points, and $b:=|G\cap H|\le r-(T-x)$, since $H$ avoids $T-x$ points of $G$.

\emph{If $b\le r/2$}, then $b\le m$, and $E,G,H$ is a good triple ($H$ misses $Y$) with all pair intersections at most $m$; Lemma~\ref{lem:L18} gives a contradiction as before.

\emph{If $b>r/2$}, all triple intersections of $E,F,G,H$ are empty and Lemma~\ref{lem:L41} applies (its gap hypothesis holds by the definition of $m$, and $m\le r/4$). Its conditions hold: $\ceil{r/2}+2m\le\ceil{r/2}+2r-2T\le T$ since $3T\ge\tfrac52r+12$; $x+m\le r-\tfrac T2+\tfrac m2+\tfrac12\le\tfrac32r-T+\tfrac12\le T$; and, using $b\le r-T+x$ and $x\le r-\tfrac{T+m}{2}+\tfrac12$,
\[
2x+b+\ceil{r/2}+2m\le 3x+r-T+\tfrac r2+\tfrac12+2m\le\tfrac{9r+m-5T+4}{2}\le3T,
\]
the last because $m\le r-T$ and $12T\ge10r+4$. Lemma~\ref{lem:L41} contradicts property $(7,2)$.
\end{proof}

\section{Proof of Theorem~\ref{thm:main}}\label{sec:main}

Let $k\ge1000$, $r=k$, $\beta=173/200$, $h=23/50$, $\ell=43/200$ (so $\ell=2-\beta-2h$), and let $\cH$ be an $r$-uniform family with property $(7,2)$. Suppose, for a contradiction, that $\tau(\cH)>K:=\ceil{\beta r}+10$. Note $K\le r$ (as $\tfrac{27}{200}r\ge11$). By Lemma~\ref{lem:pad} this also covers families of rank at most $k$.

\medskip\noindent\textbf{Step 1: no pair intersection in $[\ell r,hr]$.} Let $h_0=\floor{hr}$. Some pair intersection is at most $r-K<h_0$ (request an edge avoiding $K$ points of an edge). Let $m$ be the largest pair intersection that is at most $h_0$, attained by $E,G$. Suppose $m\ge\ell r$. Lemma~\ref{lem:balanced} with $T_0=K$ gives $F\ne E,G$ with $E\cap F\cap G=\emptyset$ and
\[
|E\cap F|,|G\cap F|\le r-\frac{K+m-1}{2}\le r-\frac{\beta r+10+\ell r-1}{2}=hr-4.5<h_0,
\]
using $\beta+\ell=2-2h$. By the maximality of $m$ both are at most $m$. Exchange the names $E,G$ if necessary so that $y:=|E\cap F|\ge z:=|G\cap F|$. Then $m\le h_0\le hr=\tfrac{23}{50}r$ and
\[
m+2y\le m+2r-(K+m-1)=2r-K+1\le(2-\beta)r-9<\tfrac{227}{200}r .
\]
Proposition~\ref{prop:local} with $T=K$ (note $\beta r+3\le K\le r$) produces a bad subfamily, contradicting property $(7,2)$. Hence $m<\ell r$. Consequently every pair intersection is either less than $\ell r$ or at least $h_0+1>hr$.

\medskip\noindent\textbf{Step 2: no pair intersection in $[hr,r/2]$.} This is Lemma~\ref{lem:gapext} (the alternative ``bad subfamily'' is excluded by property $(7,2)$).

\medskip\noindent\textbf{Step 3.} Every pair intersection is now less than $\ell r$ or greater than $r/2$. Since $\ell=\tfrac{43}{200}\le\tfrac{3\beta-2}{2}=\tfrac{119}{400}$, $\beta\ge5/6$ and $\ceil{\beta r}+4\le r$, Lemma~\ref{lem:finish} (with $\beta'=\beta$) gives $\tau(\cH)\le\ceil{\beta r}+4<K$, a contradiction. This proves Theorem~\ref{thm:main}. \qed

\begin{remark}[rounding]
The additive constant $10$ arises as follows: Step 1 needs $r-(K+m-1)/2<h_0$ and $m+2y\le(2-\beta)r$ with room for the rounding in Lemma~\ref{lem:balanced}; Lemma~\ref{lem:gapext} and Lemma~\ref{lem:finish} use $T=\ceil{\beta r}+4$; Proposition~\ref{prop:local} needs $T\ge\beta r+3$. We made no attempt to optimise the constant or the range $k\ge1000$. (All numerical inequalities of this section and of Section~\ref{sec:gap} were also checked by exact integer arithmetic for every $1000\le r\le 20000$; see Appendix~\ref{app:checks}.)
\end{remark}

\section{Further structural lemmas (independent of Theorem~\ref{thm:main})}\label{sec:further}

The results of this section are not used in the proof of Theorem~\ref{thm:main}. They are tools toward the conjecture $f(k,7)=(3/4+o(1))k$. Here $\cH$ is any family of finite nonempty sets (no uniformity) with property $(7,2)$ and $\tau(\cH)\ge t$; thus every set of at most $t-1$ points is avoided by an edge. All of them rest on the following form of the Fano construction of~\cite{FKW}.

\begin{lemma}[Fano labelling]\label{lem:fano}
Let $\Pi$ be the Fano plane. Suppose edges $G_l$ are indexed by the seven lines $l$ of $\Pi$ (repetitions allowed) and there is a map $\pi$ from the vertex set to the points of $\Pi$ such that $v\in G_l$ implies $\pi(v)\notin l$. Then $\{G_l\}$ has no transversal of size at most two.
\end{lemma}
\begin{proof}
For vertices $v,w$ (possibly equal) choose a line $l$ through $\pi(v)$ and $\pi(w)$; then $v,w\notin G_l$.
\end{proof}

\begin{theorem}[span of good triples]\label{thm:GT}
If $G_1,G_2,G_3\in\cH$ satisfy $G_1\cap G_2\cap G_3=\emptyset$, then $|G_1\cup G_2\cup G_3|\ge2t-2$. Equivalently, if $\cH$ has property $(7,2)$ and contains a good triple of union size $N$, then $\tau(\cH)\le\floor{N/2}+1$.
\end{theorem}
\begin{proof}
Let $U=G_1\cup G_2\cup G_3$, $N=|U|$, and suppose $N\le2t-3$. Label the Fano points $p_0,a,a',b,b',c,c'$ with lines $\{p_0,a,a'\}$, $\{p_0,b,b'\}$, $\{p_0,c,c'\}$ (the pencil through $p_0$) and $\{a,b,c\}$, $\{a',b',c\}$, $\{a',b,c'\}$, $\{a,b',c'\}$. Put $\mathcal A=U\setminus G_1$, $\mathcal B=G_1\setminus G_2$, $\mathcal C=G_1\cap G_2$ (a partition of $U$), and split $\mathcal A=\mathcal A_a\sqcup\mathcal A_{a'}$, $\mathcal B=\mathcal B_b\sqcup\mathcal B_{b'}$, $\mathcal C=\mathcal C_c\sqcup\mathcal C_{c'}$ as described below. Let $\pi$ map $\mathcal A_a$ to $a$, etc., and every vertex outside $U$ to $p_0$. Assign $G_1$ to $\{p_0,a,a'\}$ (it avoids $\mathcal A$ and lies in $U$), $G_2$ to $\{p_0,b,b'\}$ (it avoids $\mathcal B=G_1\setminus G_2$), and $G_3$ to $\{p_0,c,c'\}$ (it avoids $\mathcal C=G_1\cap G_2$ because the triple has no common point). For each of the four lines missing $p_0$, request an edge avoiding the union of its three classes. Then Lemma~\ref{lem:fano} applies: a vertex outside $U$ is labelled $p_0$, and the only lines through $p_0$ are the pencil lines, whose edges $G_1,G_2,G_3\subseteq U$ do not contain it; a vertex of $U$ labelled, say, $a$ lies in $\mathcal A$, hence not in $G_1$, and not in the two requested edges of the lines through $a$. We obtain a bad subfamily provided each of the four requested sets has at most $t-1$ points.

Each class occurs in exactly two of the four lines missing $p_0$, so the four loads sum to $2N$. Write $\epsilon_A=|\mathcal A_a|-|\mathcal A_{a'}|$, and similarly $\epsilon_B,\epsilon_C$; each can be chosen in $\{0\}$ (even class) or $\{\pm1\}$ (odd class). The loads equal $N/2+\tfrac12(\epsilon_A+\epsilon_B+\epsilon_C)$, $N/2+\tfrac12(-\epsilon_A-\epsilon_B+\epsilon_C)$, $N/2+\tfrac12(-\epsilon_A+\epsilon_B-\epsilon_C)$, $N/2+\tfrac12(\epsilon_A-\epsilon_B-\epsilon_C)$. With no odd class all loads are $N/2$; with one odd class they are $N/2\pm\frac12$; with two odd classes some choice gives loads in $\{N/2-1,N/2,N/2+1\}$; with three odd classes, all $\epsilon=-1$ gives loads $N/2-\tfrac32,N/2+\tfrac12,N/2+\tfrac12,N/2+\tfrac12$. In every case the maximum load is at most $\floor{N/2}+1\le t-1$ when $N\le2t-3$. So each request is legal, giving a bad subfamily of seven edges, a contradiction.
\end{proof}

\begin{remark}
Theorem~\ref{thm:GT} is sharp for the complete hypergraph $K^{(5)}_9$, which has property $(7,2)$: seven $5$-subsets of a $9$-set without a $2$-point transversal would have complements forming seven $4$-sets that cover all pairs of the $9$ points, but the covering number $C(9,4,2)$ equals $8>7$ (checked by an exact integer program; see Appendix~A). Hence then $t=\tau(K^{(5)}_9)=5$, and $\{1,2,3,4,5\}$, $\{1,2,6,7,8\}$, $\{3,4,6,7,8\}$ is a good triple of union size $8=2t-2$. The parity choice of the splits is what gives $2t-2$.
\end{remark}

\begin{theorem}[two-colour lemma]\label{thm:TC}
Let $E_0,b_1,b_2,c_1,c_2\in\cH$ with $b_1\cap b_2\cap E_0=\emptyset=c_1\cap c_2\cap E_0$, and let $E_0=D_A\sqcup D_B$ be a partition with
\[
E_0\cap\bigl((b_1\cap c_1)\cup(b_2\cap c_2)\bigr)\subseteq D_A,\qquad E_0\cap\bigl((b_1\cap c_2)\cup(b_2\cap c_1)\bigr)\subseteq D_B .
\]
Then at least one of
\[
T_A=D_A\cup(b_1\cap c_1)\cup(b_2\cap c_2),\qquad T_B=D_B\cup(b_1\cap c_2)\cup(b_2\cap c_1)
\]
is a transversal of $\cH$.
\end{theorem}
\begin{proof}
Suppose not, and let $g_1,g_2\in\cH$ avoid $T_A$, $T_B$ respectively. Label the Fano points $p_0,p_1,q,r_1,r_2,r_3,r_4$ with lines
\[
\begin{gathered}
L=\{p_0,p_1,q\},\quad B_1=\{p_0,r_1,r_2\},\quad B_2=\{p_0,r_3,r_4\},\quad C_1=\{p_1,r_1,r_3\},\\ C_2=\{p_1,r_2,r_4\},\quad M_1=\{q,r_1,r_4\},\quad M_2=\{q,r_2,r_3\},
\end{gathered}
\]
and assign $E_0,b_1,b_2,c_1,c_2,g_1,g_2$ to $L,B_1,B_2,C_1,C_2,M_1,M_2$. We define $\pi$ so that $v\in G_l\Rightarrow\pi(v)\notin l$, and apply Lemma~\ref{lem:fano}.

\emph{Vertices of $E_0$.} A vertex of $D_A$ gets $r_1$ or $r_4$ (points of $M_1$; $g_1$ misses $D_A\subseteq T_A$). The label $r_1$ is forbidden only if $v\in b_1\cup c_1$ and $r_4$ only if $v\in b_2\cup c_2$; both are forbidden only if $v\in(b_1\cap c_2)\cup(b_2\cap c_1)$ (recall $v$ is not in both $b$'s nor both $c$'s), which is impossible for $v\in D_A$. Symmetrically a vertex of $D_B$ gets $r_2$ or $r_3$ ($r_2$ forbidden by $b_1\cup c_2$, $r_3$ by $b_2\cup c_1$), which is possible since $v\notin(b_1\cap c_1)\cup(b_2\cap c_2)$, and $g_2$ misses $D_B$.

\emph{Vertices outside $E_0$.} If $v\notin b_1\cup b_2$ put $\pi(v)=p_0$ (lines $L,B_1,B_2$); if $v\notin c_1\cup c_2$ put $\pi(v)=p_1$ (lines $L,C_1,C_2$). Otherwise $v$ lies in one of $b_1\cap c_1$, $b_2\cap c_2$ (subsets of $T_A$) or $b_1\cap c_2$, $b_2\cap c_1$ (subsets of $T_B$). Take $v\in b_1\cap c_1$; then $v\notin g_1$. If $v\notin g_2$ put $\pi(v)=q$ (lines $L,M_1,M_2$). If $v\in g_2$ then $v\notin T_B$, so $v\notin c_2$ and $v\notin b_2$; the lines containing $v$'s edges are among $B_1,C_1,M_2$, whose union misses $r_4$; put $\pi(v)=r_4$. The other three cases are identical: $v\in b_2\cap c_2$, $v\in g_2$ gives lines among $B_2,C_2,M_2$, missing $r_1$; $v\in b_1\cap c_2$, $v\in g_1$ gives lines among $B_1,C_2,M_1$, missing $r_3$; $v\in b_2\cap c_1$, $v\in g_1$ gives lines among $B_2,C_1,M_1$, missing $r_2$.

By Lemma~\ref{lem:fano} the seven edges form a bad subfamily, contradicting property $(7,2)$.
\end{proof}

\begin{theorem}[$K_4$ criterion]\label{thm:K4}
Let $E_0\in\cH$. Suppose $E_0=E_1\sqcup E_2\sqcup E_3\sqcup E_4$ and there are six edges $G_{ab}$ ($1\le a<b\le4$) with $G_{ab}\cap E_0\subseteq E_a\cup E_b$ such that no vertex outside $E_0$ lies in $G_{ab}\cap G_{bc}\cap G_{ac}$ for any three distinct $a,b,c$. Then $\{E_0\}\cup\{G_{ab}\}$ is a bad subfamily. Conversely, every seven edges indexed by Fano lines and admitting a labelling as in Lemma~\ref{lem:fano}, one of which is $E_0$, arise in this way.
\end{theorem}
\begin{proof}
Take a line $L=\{s_1,s_2,s_3\}$ and the four points $r_1,\dots,r_4$ off $L$. Each of the six other lines contains exactly one $s_i$ and two points $r_c,r_d$; assign to it the edge $G_{ab}$, $\{a,b\}=[4]\setminus\{c,d\}$, and assign $E_0$ to $L$. Label $v\in E_a$ by $r_a$: the lines through $r_a$ other than $L$ carry $G_{bc},G_{bd},G_{cd}$ with $\{b,c,d\}=[4]\setminus\{a\}$, whose traces on $E_0$ avoid $E_a$. For $v\notin E_0$ let $\sigma(v)$ be the set of pairs $ab$ with $v\in G_{ab}$, viewed as edges of $K_4$. The lines through $s_i$ carry the edges of a perfect matching of $K_4$, and the lines through $r_a$ (other than $L$) carry the triangle on $[4]\setminus\{a\}$. So $v$ can be labelled iff $\sigma(v)$ is disjoint from some perfect matching or from some triangle. If $\sigma(v)$ meets all three perfect matchings, it contains three edges, one from each matching, and such a triple is either a triangle or a star; a triangle-free $\sigma(v)$ meeting all matchings contains a star $\{ab,ac,ad\}$ and no further edge (adding any edge creates a triangle), hence it is disjoint from the triangle on $\{b,c,d\}$. Thus $v$ can be labelled whenever $\sigma(v)$ contains no triangle, which is the hypothesis. Lemma~\ref{lem:fano} gives the first claim. Conversely, given such a labelled Fano tuple with $E_0$ on $L$, the vertices of $E_0$ carry labels off $L$; let $E_a$ be those labelled $r_a$, and the same dictionary shows the trace condition, while a vertex outside $E_0$ lying in a triangle $G_{ab}\cap G_{bc}\cap G_{ac}$ would meet all matchings and all triangles and could not be labelled.
\end{proof}

\begin{theorem}[spread lemma]\label{thm:spread}
Let $E_0\in\cH$ and $E_0=E_1\sqcup E_2\sqcup E_3\sqcup E_4$, and put $s=t-1-\max(|E_1\cup E_4|,|E_2\cup E_3|)\ge0$. Let $b_1,b_2\in\cH$ with $b_1\cap E_0\subseteq E_3\cup E_4$ and $b_2\cap E_0\subseteq E_1\cup E_2$. Let $\mu$ and $\nu$ be probability distributions on $\mathcal O_{24}=\{G\in\cH:G\cap E_0\subseteq E_2\cup E_4\}$ and $\mathcal O_{13}=\{G\in\cH:G\cap E_0\subseteq E_1\cup E_3\}$, and let $p=\max_{v\notin E_0}\mu(G\ni v)$, $q=\max_{v\notin E_0}\nu(G\ni v)$. Then
\[
|(b_1\cup b_2)\setminus E_0|\,(p+q)\ \ge\ s+1 .
\]
In particular, if $\cH$ has rank at most $k$ and the partition is balanced ($|E_1\cup E_4|=\ceil{|E_0|/2}$, $|E_2\cup E_3|=\floor{|E_0|/2}$), then $t\le\ceil{|E_0|/2}+2k(p+q)$; if moreover $p,q\le1/16$ then $t\le\ceil{k/2}+k/4$.
\end{theorem}
\begin{proof}
For $c_1\in\mathcal O_{24}$ and $c_2\in\mathcal O_{13}$ apply Theorem~\ref{thm:TC} with $D_A=E_1\cup E_4$, $D_B=E_2\cup E_3$: the hypotheses hold since $b_1\cap b_2\cap E_0=\emptyset=c_1\cap c_2\cap E_0$, $b_1\cap c_1\cap E_0\subseteq E_4$, $b_2\cap c_2\cap E_0\subseteq E_1$, $b_1\cap c_2\cap E_0\subseteq E_3$ and $b_2\cap c_1\cap E_0\subseteq E_2$. So $T_A$ or $T_B$ is a transversal and has at least $t$ points. Both are contained in $D_A\cup W$, resp.\ $D_B\cup W$, where $W=((b_1\cup b_2)\cap(c_1\cup c_2))\setminus E_0$; hence $|W|\ge s+1$ for every such $c_1,c_2$. Draw $c_1\sim\mu$ and $c_2\sim\nu$ independently. Then
\[
s+1\le\mathbb E|W|\le\sum_{v\in(b_1\cup b_2)\setminus E_0}\bigl(\Pr[v\in c_1]+\Pr[v\in c_2]\bigr)\le|(b_1\cup b_2)\setminus E_0|(p+q).
\]
For the consequences use $|(b_1\cup b_2)\setminus E_0|\le2k$ and $s+1=t-\ceil{|E_0|/2}$ (if $s<0$ the bound $t\le\ceil{|E_0|/2}$ is immediate).
\end{proof}

\begin{remark}
Theorems~\ref{thm:TC}--\ref{thm:spread} only produce bad subfamilies of \emph{Fano type} (Lemma~\ref{lem:fano}). By an example in the working note (two disjoint parts, coefficient approaching $4/5$ in a continuous type model), Fano-type tuples alone cannot handle families with disjoint edges; any proof of the $3/4$ bound along these lines must also use other bad seven-tuples. Theorem~\ref{thm:spread} shows that in a hypothetical family with $t\ge(3/4+\varepsilon)k$ and rank $k$, for every edge and every balanced quartering, the two ``outside clouds'' cannot both be $\tfrac1{16}$-spread.
\end{remark}

\section{A threshold theorem for symmetric families}\label{sec:threshold}

The extremal example for the conjectured value $3/4$, the complete hypergraph $K^{(k)}_{\lfloor 7k/4\rfloor-1}$, is invariant under all permutations of its ground set, and the parity example of~\cite{FKW} is invariant under all permutations preserving a fixed set $X$. It is therefore natural to test the conjecture on families with a large symmetry group. In this section we prove it, up to an additive constant, for a natural class of such families: unions of \emph{threshold} conditions.

\begin{definition}\label{def:threshold}
Let $V=P_1\sqcup\dots\sqcup P_p$ be a partition, let $I\subseteq[p]$, and let $t_i\ (i\in I)$ be integers with $1\le t_i\le k$. The \emph{threshold family} $\mathcal T=\mathcal T(P;I;t)$ consists of all $k$-subsets $E\subseteq V$ such that $|E\cap P_i|\ge t_i$ for at least one $i\in I$. We call $P_i$ ($i\in I$) the \emph{boxes}.
\end{definition}

A threshold family is invariant under $\mathrm{Sym}(P_1)\times\dots\times\mathrm{Sym}(P_p)$. The complete hypergraph is the case $p=1$, $t_1\le k$.

\begin{theorem}\label{thm:threshold}
Let $\mathcal T$ be a threshold family with at least three boxes, and suppose $|P_i|\ge t_i+7$ for every box $P_i$. If $\mathcal T$ has property $(7,2)$, then
\[
\tau(\mathcal T)\le \tfrac34k+8p .
\]
\end{theorem}

The bad subfamilies we construct are always of the Fano type used in Section~\ref{sec:further}. We first isolate the only fact about them that we need.

\begin{lemma}[Fano windows]\label{lem:fanowindows}
Label the points of a Fano plane by $\{1,\dots,7\}$. Let $C_1,\dots,C_7$ be pairwise disjoint sets, and for every line $\ell$ let $G_\ell$ be a set contained in the \emph{window} $W_\ell=\bigcup_{q\notin\ell}C_q$. Then no two points meet all seven sets $G_\ell$.
\end{lemma}

\begin{proof}
A point of $C_q$ lies only in sets $G_\ell$ with $q\notin\ell$. Given $u\in C_q$ and $v\in C_{q'}$ (possibly $q=q'$), choose a line $\ell$ through $q$ and $q'$; then $G_\ell$ misses both. A point outside $\bigcup C_q$ lies in no $G_\ell$.
\end{proof}

\subsection{The continuous model}
Scale by $k$. Put $x_i=|P_i|/k$ (the \emph{capacities}), $X=\sum_ix_i$ and $\theta_i=t_i/k$. A \emph{type} is a vector $a\in\mathbb R^p_{\ge0}$ with $a\le x$, $\sum_ia_i=1$ and $a_i\ge\theta_i$ for some box $i$; types are the rescaled profiles $(|E\cap P_i|/k)_i$ of edges. A vector $0\le u\le x$ is \emph{free} if no type satisfies $a\le u$, and
\[
\tau^*=X-\sup\Bigl\{\textstyle\sum_iu_i: u\text{ free}\Bigr\}.
\]

\begin{lemma}\label{lem:tauint}
$\tau(\mathcal T)\le k\,\tau^*+p$.
\end{lemma}

\begin{proof}
Let $u$ be free and let $S\subseteq V$ have $\lfloor ku_i\rfloor$ points in each $P_i$. An edge $E\subseteq S$ would give a type $|E\cap P|/k\le u$; so $S$ is independent, and $\tau(\mathcal T)\le|V|-|S|\le kX-k\sum u_i+p$. Take the supremum over $u$.
\end{proof}

We say a box $i$ is \emph{effective} if some type $a$ has $a_i\ge\theta_i$. For an effective box put $\theta_i'=\max(\theta_i,\,1-X+x_i)$; every type with $a_i\ge\theta_i$ has $a_i\ge\theta'_i$, because the remaining coordinates sum to at most $X-x_i$. Put $d_i=x_i-\theta_i'\ge0$.

\begin{lemma}\label{lem:freevector}
$\tau^*\le X-1$, and $\tau^*\le\sum_{i\text{ effective}}d_i$.
\end{lemma}

\begin{proof}
Every $u$ with $\sum u_i<1$ is free. The vector with coordinates $\theta'_i-\varepsilon$ on effective boxes and $x_j$ elsewhere is free, since a type below it would have $a_i<\theta'_i$ for every effective box $i$ and could not reach any threshold. Let $\varepsilon\to0$.
\end{proof}

A \emph{Fano placement} for the template $T(A,B,C)$, where $A,B,C$ are distinct effective boxes, is the following continuous object. Fix a Fano point $p_0$; call the four lines missing $p_0$ the \emph{quadrilateral} and the three lines through $p_0$ the \emph{pencil} $\ell_1,\ell_2,\ell_3$. Assign box $A$ to the quadrilateral lines, box $B$ to $\ell_1,\ell_2$ and box $C$ to $\ell_3$, and write $\beta(\ell)$ for the box of line $\ell$. A placement consists of numbers $c_{i,q}\ge0$ (the mass of Fano point $q$ in part $i$) and $a^\ell_i\ge0$ (the trace of row $\ell$ in part $i$) with
\begin{equation}\label{eq:placement}
\sum_qc_{i,q}\le x_i,\qquad a^\ell_i\le\sum_{q\notin\ell}c_{i,q},\qquad\sum_ia^\ell_i=1,\qquad a^\ell_{\beta(\ell)}\ge\theta'_{\beta(\ell)} .
\end{equation}

\begin{lemma}[realisation]\label{lem:realise}
Delete $\min(7,|P_i|)$ points from every part, obtaining the threshold family $\mathcal T'$ on the smaller parts. If $\mathcal T'$ admits a Fano placement, or a type with $a\le\tfrac47x$, then $\mathcal T$ does not have property $(7,2)$.
\end{lemma}

\begin{proof}
Take the continuous data for $\mathcal T'$, with capacities $x'_i\ge x_i-7/k$. Let $C_{i,q}$ be disjoint subsets of $P_i$ of sizes $\lceil kc_{i,q}\rceil$; this fits, since $\sum_q\lceil kc_{i,q}\rceil\le kx'_i+7\le|P_i|$. The window of row $\ell$ then has at least $\lceil ka^\ell_i\rceil$ points in $P_i$. Choose $b^\ell_i=\lceil ka^\ell_i\rceil$ points there. This gives at least $k$ points with at least $\lceil ka^\ell_{\beta}\rceil\ge t_\beta$ of them in the box $\beta=\beta(\ell)$. Discard points outside the box first, and if necessary box points down to $t_\beta$, to get exactly $k$ points. The result is an edge $G_\ell$ of $\mathcal T$ inside its window, and by Lemma~\ref{lem:fanowindows} the seven edges have no $2$-point transversal. For a type $a\le\frac47x'$ use seven classes of size $\lceil kx'_i/7\rceil$ in every part and seven copies of that type.
\end{proof}

\subsection{The continuous theorem}

\begin{theorem}\label{thm:threshold-cont}
In the continuous model, suppose there are at least three effective boxes and $\tau^*\ge 3/4$. Then either some type satisfies $a\le\frac47x$, or for any three effective boxes $A,B,C$ the template $T(A,B,C)$ admits a Fano placement.
\end{theorem}

\begin{proof}
\emph{Step 1 (the homogeneous case).} By Lemma~\ref{lem:freevector}, $X\ge7/4$, so $\sum_i\frac47x_i\ge1$. If some effective box has $\theta'_i\le\frac47x_i$, the vector with $a_i=\theta'_i$, filled up to total $1$ inside $\frac47x$, is a type below $\frac47x$. Assume from now on that $\frac47x_i<\theta'_i\le\min(x_i,1)$ for every effective box.

\emph{Step 2 (one light part).} Rows of the template only need thresholds in $A,B,C$. Merge all other parts into one part $L$ of capacity $x_L$. The constraints~\eqref{eq:placement} are linear and homogeneous in $(c_{i,\cdot},a_{i}^{\cdot},x_i)$, so a placement for the merged part splits back into the original parts proportionally to their capacities. By Step~1, each merged effective box has $d_j<\frac37x_j$, and non-boxes contribute nothing to Lemma~\ref{lem:freevector}. Writing $d_X=x_X-\theta'_X$ for $X\in\{A,B,C\}$, we get
\begin{equation}\label{eq:g}
g:=d_A+d_B+d_C+\tfrac37x_L\ \ge\ \tau^*\ \ge\ \tfrac34 .
\end{equation}
If $x_L\ge7/4$, give each row trace $\theta'_{\beta(\ell)}$ in its own box and the rest in $L$. Put mass $\theta'_A$ on $p_0$ in part $A$, and mass $\theta'_B$ (resp.\ $\theta'_C$) on a point of part $B$ (resp.\ $C$) off $\ell_1\cup\ell_2$ (resp.\ off $\ell_3$). Put $x_L/7\ge\frac14$ on every point in $L$. This is a placement, so assume $x_L\le7/4$.

\emph{Step 3 (convexity).} For fixed $(x_A,x_B,x_C,\theta'_A,\theta'_B,\theta'_C,x_L)$, conditions~\eqref{eq:placement} are linear inequalities in the variables $c,a$ whose right-hand sides depend linearly on these parameters. The set of parameters admitting a placement is therefore the projection of a polyhedron, and is convex. It suffices to show that it contains every vertex of the domain
\[
D=\{(x_X,\theta'_X)_{X=A,B,C}\in\Delta^3,\ 0\le x_L\le\tfrac74,\ g\ge\tfrac34\},\qquad
\Delta=\{(x,\theta):\tfrac47x\le\theta\le\min(x,1)\}.
\]

\emph{Step 4 (the vertices are degenerate).} $\Delta$ is the triangle with vertices $(0,0)$, $(1,1)$ and $(\frac74,1)$, on which $d=x-\theta$ takes the values $0,0,\frac34$. At the vertices of the product $P=\Delta^3\times[0,\frac74]$ the function $g$ therefore takes values in $\frac34\mathbb Z$. Along an edge of $P$ only one factor moves, so $g$ changes by $0$ or $\frac34$, and the hyperplane $g=\frac34$ can meet an edge of $P$ only at an endpoint. Hence every vertex of $D=P\cap\{g\ge\frac34\}$ is a vertex of $P$. At such a vertex each of $A,B,C$ is \emph{empty} ($x=\theta=0$), \emph{tight} ($x=\theta=1$) or a \emph{Fano part} ($x=\frac74$, $\theta=1$), and $x_L\in\{0,\frac74\}$. Since $g\ge\frac34$, at least one part of capacity $\frac74$ exists (a Fano part, or $L$); call it the \emph{host}.

\emph{Step 5 (the vertices).} Give every row its whole mass $1$ in its own box if that box is tight or a Fano part, and in the host if its box is empty; the threshold conditions then hold ($\theta'=1$ or $\theta'=0$). A tight part $A$ carries the four quadrilateral rows: put mass $1$ on $p_0$, which lies in the window of every line missing $p_0$. A tight $B$ carries the rows $\ell_1,\ell_2$: put mass $1$ on one of the two points off $\ell_1\cup\ell_2$. A tight $C$ carries $\ell_3$: put mass $1$ on a point off $\ell_3$. Every part of capacity $\frac74$ gets mass $\frac14$ on each Fano point; each of its windows then has mass $1$ and can carry any row. This is a placement at every vertex of $D$. By Step~3 it exists on all of $D$, and the parameters of our family lie in $D$ by Steps~1 and~2.
\end{proof}

\begin{proof}[Proof of Theorem~\ref{thm:threshold}]
Suppose $\tau(\mathcal T)>\frac34k+8p$ and delete $\min(7,|P_i|)$ points from each part to obtain $\mathcal T'$. Then $\tau(\mathcal T')\ge\tau(\mathcal T)-7p>\frac34k+p$. Every box keeps $|P'_i|\ge t_i$ points, so it stays effective when $X'\ge1$, and $\mathcal T'$ is nonempty. By Lemma~\ref{lem:tauint}, $\tau^*(\mathcal T')>\frac34$. Theorem~\ref{thm:threshold-cont} and Lemma~\ref{lem:realise} show that $\mathcal T$ is not $(7,2)$.
\end{proof}

\begin{remark}\label{rem:threshold}
(i) The constant $\frac34$ is sharp even for $p=1$, by the complete hypergraph.
(ii) At least three boxes are needed for Fano-type subfamilies. With two boxes of capacity $\frac{11}8$ and thresholds equal to $1$, no Fano placement exists although $\tau^*=\frac34$. That family is nevertheless not $(7,2)$: four edges inside one box with no common point, together with one edge inside the other box, have no $2$-point transversal. So the two-box case needs other bad subfamilies.
\end{remark}


\section{Concluding remarks}

\emph{The general bound.} The chain of Section~\ref{sec:main} cannot be pushed much further with the same ingredients. A numerical optimisation over four-request static templates (not a proof) suggests that the best static budget at the binding configuration of Proposition~\ref{prop:local} is about $0.862k$, so new adaptive ideas are needed below roughly $0.862$.

\emph{Symmetric families.} Theorem~\ref{thm:threshold} is one instance of the following question, which is equivalent to Erd\H{o}s's question restricted to families invariant under $\prod_i\mathrm{Sym}(P_i)$ for a partition into boundedly many parts:

\smallskip
\noindent\textbf{Question.} \emph{Is $\tau(\mathcal H)\le\frac34k+O(p)$ for every $(7,2)$-family $\mathcal H$ of $k$-sets that is invariant under $\mathrm{Sym}(P_1)\times\dots\times\mathrm{Sym}(P_p)$?}
\smallskip

The method of Section~\ref{sec:threshold} works for other classes with a polyhedral parameter space. It has limits, though: Fano-type bad subfamilies alone do not suffice in general. For two parts of capacity $x\in(\frac98,\frac97)$ and the two types $(s,1-s)$ and $(1-s,s)$ with $s<1-\frac{2x}3$, the continuous transversal coefficient tends to $\frac67$ as $x\to\frac97$. Yet an exact rational certificate shows that none of the $2^7$ ways of assigning the two types to the seven Fano rows is realisable. Such families are nevertheless not $(7,2)$, because a non-Fano bad $7$-tuple exists. Any proof of the $3/4$ bound must therefore use bad subfamilies of more than one shape.

\appendix
\section{Independent computer checks (sanity supplement)}\label{app:checks}

The proof above does not depend on any computation. The following exact checks were run as independent confirmation; all use exact rational or integer arithmetic.
\begin{enumerate}
\item \texttt{w5\_paper\_prop10\_vertices.py} (written for this manuscript, independent of earlier scripts): enumerates the vertices of the closure of each of the nine case polytopes of Proposition~\ref{prop:local} (a, b1, b2, b3, c1, c2, c3, c4 split A, c4 split B) and checks, at every vertex, every lemma hypothesis as stated in Section~\ref{sec:closing} and every auxiliary inequality asserted in the text, including (C), $(*)$ and the condition (ii) of case (c2). Since each hypothesis is a convex piecewise-linear function bounded by a constant (or affine on each split branch), vertex checks cover the polytopes. Result: all pass; every case polytope is nonempty (a: 4, b1: 14, b2: 10, b3: 4, c1: 6, c2: 5, c3: 6, c4A: 6, c4B: 6 vertices); a randomised exact test finds no point of the region outside the union of the case polytopes. A mutation (budget $0.8645$) produces failures.
\item \texttt{w5\_paper\_constants.py}: checks with exact integers every numerical inequality of Section~\ref{sec:main}, Lemma~\ref{lem:gapext}, Lemma~\ref{lem:finish} (at $\beta'=\beta$) and the budget of Lemma~\ref{lem:L18} at $m=\floor{119r/400}$, for every $r$ with $1000\le r\le20000$ (and in fact $80\le r\le 20000$). Result: all pass.
\item Referee scripts from an earlier review round (same directory): \texttt{referee\_w4\_handbound\_vertices.py} (vertex enumeration of the case polytopes), \texttt{referee\_w4\_handbound\_fm.py} (Fourier--Motzkin elimination with strictness flags, showing each case system together with the negation of each lemma inequality is infeasible), \texttt{referee\_w4\_handbound\_int.py} (exhaustive integer check of Proposition~\ref{prop:local} with $T=\ceil{\beta r}+3$ and the rounded S1 splits: $9{,}980{,}329$ triples at $r=1000$ and all triples at several smaller $r$, no failure), \texttt{referee\_w4\_handbound\_indep.py}, \texttt{referee\_w4\_handbound\_s1s2.py}, \texttt{referee\_w4\_handbound\_static\_indep.py} (explicit set systems: every candidate pair of S1 and S2 is covered and all request sizes are at most $T$; the S1 closed form agrees with the exact linear-programming optimum), and \texttt{referee\_w4\_handbound\_exh.py} (for all $r\le12$, an exhaustive first-response adversary for Lemmas~\ref{lem:L26}, \ref{lem:L32}, \ref{lem:L31} with exact treatment of the final responses; no failure, while the mutation $T-1$ fails).
\end{enumerate}
For this manuscript we re-ran \texttt{referee\_w4\_handbound\_vertices.py}, \texttt{referee\_w4\_handbound\_fm.py} and \texttt{referee\_w4\_handbound\_int.py} (the latter at $r=200,333,401$ and exhaustively at $r=1000$, $9{,}980{,}329$ triples); all pass. The remaining scripts of item 3 were run in the earlier review round and were not re-run here. All these checks confirm arithmetic and case coverage only; the mathematical content is in the proofs.

\begin{thebibliography}{9}
\bibitem{BKS} M.~Buci\'c, D.~Kor\'andi and B.~Sudakov, \emph{Covering graphs by monochromatic trees and Helly-type results for hypergraphs}, Combinatorica \textbf{41} (2021), 319--352; arXiv:1902.05055.
\bibitem{EFKT} P.~Erd\H{o}s, D.~G.~Fon-Der-Flaass, A.~V.~Kostochka and Zs.~Tuza, \emph{Small transversals in uniform hypergraphs}, Siberian Adv.\ Math.\ \textbf{2} (1992), 82--88.
\bibitem{EHT} P.~Erd\H{o}s, A.~Hajnal and Zs.~Tuza, \emph{Local constraints ensuring small representing sets}, J.\ Combin.\ Theory Ser.~A \textbf{58} (1991), 78--84.
\bibitem{EP} Erd\H{o}s Problems website, Problem \#644, \texttt{https://www.erdosproblems.com/644} (status: open; accessed 24 September 2026; lists $f(k,3)=2k$, $f(k,4)=\lfloor3k/2\rfloor$, $f(k,5)=\lfloor5k/4\rfloor$, $f(k,6)=k$ from~\cite{EFKT}).
\bibitem{FKW} D.~G.~Fon-Der-Flaass, A.~V.~Kostochka and D.~R.~Woodall, \emph{Transversals in uniform hypergraphs with property $(7,2)$}, Discrete Math.\ \textbf{207} (1999), 277--284.
\bibitem{Kos} A.~Kostochka, \emph{Transversals in uniform hypergraphs with property $(p,2)$}, Combinatorica \textbf{22} (2002), 275--285 (cited in~\cite{BKS}; content not consulted) [citation to check].
\end{thebibliography}

\end{document}
